\documentclass[12pt]{article}

% ================== CONFIGURACIÓN GENERAL ==================
\usepackage[utf8]{inputenc}
\usepackage[T1]{fontenc}
\usepackage[spanish]{babel}
\usepackage[letterpaper, left=3cm, right=2.5cm, top=2.5cm, bottom=2.5cm]{geometry}

% ================== FORMATO DE TEXTO ==================
\usepackage{setspace}
\setstretch{1.5}                % interlineado 1.5
\usepackage{newtxtext,newtxmath} % Times New Roman (texto y matemáticas)
\usepackage{titlesec}           % control de formato de títulos
\usepackage{lipsum}             % texto de ejemplo
\usepackage{xcolor}             % colores (opcional)

% ================== MATEMÁTICAS ==================
\usepackage{amsmath, mathtools, bm} % evita conflicto con amssymb y newtxmath
\DeclareUnicodeCharacter{2265}{\ensuremath{\geq}} % ≥
\DeclareUnicodeCharacter{2264}{\ensuremath{\leq}} % ≤

% ================== TABLAS ==================
\usepackage{tabularx}   % tablas que se ajustan al ancho
\usepackage{booktabs}   % líneas finas de tabla
\usepackage{array}      % columnas tipo p, m, b
\usepackage{multirow}   % combinar filas
\usepackage{longtable}  % tablas que se parten entre páginas
\newcommand{\cbox}{\(\square\)} % casilla vacía

% ================== LISTAS Y ENUMERACIONES ==================
\usepackage{enumitem}
\setlist[enumerate]{listparindent=\parindent, parsep=0pt}

% ================== IMÁGENES ==================
\usepackage{graphicx}
\graphicspath{{./Imagenes/}}
\usepackage{float}      % para usar [H] en figuras

% ================== PDF Y REFERENCIAS ==================
\usepackage[hidelinks]{hyperref} % hipervínculos sin color
\usepackage{pdfpages}            % para incluir PDFs externos
\usepackage{comment}             % para comentar bloques de texto

% ================== CONTADORES PERSONALIZADOS ==================
\newcounter{intro}
\renewcommand{\theintro}{\arabic{intro}}

% (opcional) espacio entre secciones
% \titlespacing*{\section}{0pt}{3ex plus 1ex minus .2ex}{2ex}


%%%%%%%%%%% NECESARIO%%%%%%%%%%%%%%%%

\usepackage{enumitem}
\setlist[enumerate]{leftmargin=*, itemsep=2pt}
\setlist[itemize]{leftmargin=*, itemsep=2pt}




\begin{comment}
	
	% --- Configuración de enumitem para numeración multinivel ---
	\setlist[enumerate,1]{label=\arabic*., leftmargin=1.2cm}
	\setlist[enumerate,2]{label=\arabic{enumi}.\arabic*., leftmargin=1.4cm}
	\setlist[enumerate,3]{label=\arabic{enumi}.\arabic{enumii}.\arabic*., leftmargin=1.6cm}
	\setlist[enumerate,4]{label=\alph*) , leftmargin=1.8cm}
	\setlist[itemize]{leftmargin=1.6cm}
	
\end{comment}


\begin{comment}
	PARA INDICES CREO
	\tableofcontents
	\newpage
\end{comment}



\begin{document}
	
	%Para no generar numeracion de paginas
	\pagestyle{empty}
	
%	\thispagestyle{fancy} 
	
	%es para emparejar el tamaño de todas las letras de titulo y contenido
%	\titleformat{\section}{\large\bfseries}{\thesection.}{0.5em}{}
%	\titleformat{\subsection}{\normalsize\bfseries}{\thesubsection.}{0.5em}{}
	
	%\title{\vspace{-2cm}Área de Métodos Constructivos
%	\\ Descripción y proceso constructivo de losas}
%	\author{}
%	\date{}
	

	
	
	
%	\maketitle
	
	%\vspace{-2.5cm}
	
%
\includepdf{portadaU2/caratulaU}
%\includepdf[pages=-,pagecommand={}]{hoja/hojaadjunta}

%%-----------TITULOSSSS--------------------------------------
		\begin{center}
		\textbf{\LARGE  Métodos Constructivos}
	\end{center}

%--------------SUBTITULOS---------------------

	% Primera introducción
\stepcounter{intro}
\section*{\centering \theintro\ Proceso constructivo de elementos de concreto}
\addcontentsline{toc}{section}{\theintro\ Descripción y proceso constructivo de elementos de concreto}
\setcounter{section}{0}


 
 
 \section{Losa tradicional (maciza de concreto reforzado)}
 
 \subsection*{Descripción general}
 La \textbf{losa tradicional} es una placa maciza de concreto reforzado colada \emph{in situ} sobre cimbra. Su rigidez proviene del peralte de la losa y del acero de refuerzo en una o dos direcciones. Es un sistema versátil que se adapta a múltiples geometrías y cargas, a costa de mayor peso propio y demanda de cimbra y curado.
 
 
 
 % =====================================================================
 \subsection{Checklist de supervisión}
 
 \subsubsection*{1. Antes de iniciar (preparación)}
 \begin{tabularx}{\linewidth}{>{\raggedright\arraybackslash}p{0.93\linewidth} >{\centering\arraybackslash}m{0.07\linewidth}}
 	\toprule
 	\textbf{Actividad a revisar / ordenar} & \textbf{Estado}\\
 	\midrule
 	Planos y especificaciones en sitio (peralte, recubrimientos, diámetros, f'c, juntas, aditivos). & \cbox \\
 	Trazos y \textbf{niveles} de losa marcados en vigas/muros; ejes verificados. & \cbox \\
 	Materiales disponibles y protegidos (acero limpio, madera, separadores, vibradores, mangueras). & \cbox \\
 	Programa de colado (volumen, acceso, bomba, personal, condiciones de clima). & \cbox \\
 	\textit{Orden a cuadrilla:} limpieza del frente y habilitar pasarelas. & \cbox \\
 	\bottomrule
 \end{tabularx}
 
 \subsubsection*{2. Encofrado (cimbra y apuntalamiento)}
 \begin{tabularx}{\linewidth}{>{\raggedright\arraybackslash}p{0.93\linewidth} >{\centering\arraybackslash}m{0.07\linewidth}}
 	\toprule
 	\textbf{Actividad a revisar / ordenar} & \textbf{Estado}\\
 	\midrule
 	Cimbra plana, alineada y \textbf{estanca} (sin fugas); bases firmes. & \cbox \\
 	Apuntalamientos con arriostres; contraflecha si la indicó el diseñador. & \cbox \\
 	Aplicación uniforme de \textbf{desmoldante} (evitar contaminar acero). & \cbox \\
 	Sellado de juntas de cimbra; pasarelas seguras para el colado. & \cbox \\
 	\textit{Orden:} reapretar uniones, calzar puntales y corregir holguras. & \cbox \\
 	\bottomrule
 \end{tabularx}
 
 \subsubsection*{3. Armado de acero}
 \begin{tabularx}{\linewidth}{>{\raggedright\arraybackslash}p{0.93\linewidth} >{\centering\arraybackslash}m{0.07\linewidth}}
 	\toprule
 	\textbf{Actividad a revisar / ordenar} & \textbf{Estado}\\
 	\midrule
 	Diámetros, separación y \textbf{recubrimientos} conforme a planos; separadores suficientes. & \cbox \\
 	Traslapes/anclajes con longitudes de desarrollo correctas; amarre con alambre recocido. & \cbox \\
 	Refuerzo adicional en bordes, sobre apoyos y alrededor de huecos/columnas. & \cbox \\
 	Malla superior (temperatura) a \textbf{cota} correcta; acero limpio, sin grasa/óxido suelto. & \cbox \\
 	Coordinación MEP: pasos de tuberías/ductos \textbf{sin interferir} al refuerzo. & \cbox \\
 	\bottomrule
 \end{tabularx}
 
 \subsubsection*{4. Pre-colado}
 \begin{tabularx}{\linewidth}{>{\raggedright\arraybackslash}p{0.93\linewidth} >{\centering\arraybackslash}m{0.07\linewidth}}
 	\toprule
 	\textbf{Actividad a revisar / ordenar} & \textbf{Estado}\\
 	\midrule
 	Cimbra, acero e instalaciones \textbf{limpios y secos}; retirar residuos/astillas. & \cbox \\
 	Vibradores probados; energía/combustible disponible; bomba lista. & \cbox \\
 	Verificar \textbf{revenimiento (slump)} y temperatura de la primera olla. & \cbox \\
 	Humedecer cimbra (sin charcos) minutos antes del vaciado. & \cbox \\
 	Pasarelas y cubetas listas para cantos y rincones. & \cbox \\
 	\bottomrule
 \end{tabularx}
 
 \subsubsection*{5. Colocación, vibrado y acabado}
 \begin{tabularx}{\linewidth}{>{\raggedright\arraybackslash}p{0.93\linewidth} >{\centering\arraybackslash}m{0.07\linewidth}}
 	\toprule
 	\textbf{Actividad a revisar / ordenar} & \textbf{Estado}\\
 	\midrule
 	Altura de descarga controlada; avance por \textbf{paños} (\(\leq\) 15\,cm). & \cbox \\
 	\textbf{Vibrado interno} adecuado (penetración rápida, retiro lento, traslape de zonas). & \cbox \\
 	Nivelación con regla/guía láser; acabado según especificación. & \cbox \\
 	Evitar acumulación de agua de exudación antes del acabado final. & \cbox \\
 	Prohibir pisar el acero sin pasarelas durante el colado. & \cbox \\
 	\bottomrule
 \end{tabularx}
 
 \subsubsection*{6. Juntas (si aplica)}
 \begin{tabularx}{\linewidth}{>{\raggedright\arraybackslash}p{0.93\linewidth} >{\centering\arraybackslash}m{0.07\linewidth}}
 	\toprule
 	\textbf{Actividad a revisar / ordenar} & \textbf{Estado}\\
 	\midrule
 	Ubicación de \textbf{juntas de construcción}; preparación de sustrato (rugoso/limpio). & \cbox \\
 	Lechada de unión/epóxicos cuando lo indique el diseño; protección de bordes. & \cbox \\
 	Cortes tempranos de contracción (si fueron especificados) en tiempo y profundidad. & \cbox \\
 	\bottomrule
 \end{tabularx}
 
 \subsubsection*{7. Curado}
 \begin{tabularx}{\linewidth}{>{\raggedright\arraybackslash}p{0.93\linewidth} >{\centering\arraybackslash}m{0.07\linewidth}}
 	\toprule
 	\textbf{Actividad a revisar / ordenar} & \textbf{Estado}\\
 	\midrule
 	Inicio inmediato del curado tras acabado final: \textbf{membrana, riego o mantas}. & \cbox \\
 	Mantener curado continuo (típ. \(\geq\) 7 días) o según especificación/ensayes. & \cbox \\
 	Proteger contra viento/sol/lluvia; evitar tránsito o acopio sobre la losa. & \cbox \\
 	\bottomrule
 \end{tabularx}
 
 \subsubsection*{8. Descimbrado y apuntalamientos}
 \begin{tabularx}{\linewidth}{>{\raggedright\arraybackslash}p{0.93\linewidth} >{\centering\arraybackslash}m{0.07\linewidth}}
 	\toprule
 	\textbf{Actividad a revisar / ordenar} & \textbf{Estado}\\
 	\midrule
 	Retiro de moldes laterales cuando la superficie no se dañe. & \cbox \\
 	\textbf{Apuntalamientos}: retirar solo con autorización y resistencia alcanzada (ensayes). & \cbox \\
 	Verificación de superficie, fisuras, flechas y nivelación final. & \cbox \\
 	\bottomrule
 \end{tabularx}
 
 
 
 
 
 %%%%%%%%%%%%%%%%%%%%%%%%%%%%%%%%%%%%%%%%%%%%%%%%%%%
 
 
 
 
 \section{Vigueta y bovedilla}
 
 \subsection*{Descripción general}
 El sistema de \textbf{vigueta y bovedilla} está formado por viguetas prefabricadas (pretensadas o de alma abierta) que trabajan como elementos portantes, y piezas de aligeramiento (bovedillas de concreto aligerado, cerámica o poliestireno) que conforman el entrevigado. Sobre el conjunto se vacía una \textbf{capa de compresión} de concreto que integra el diafragma. Requiere modulación estricta, apuntalamiento temporal y seguir el manual del proveedor.
 

 
 % =====================================================================
 \subsection{Checklist de supervisión}
 
 \subsubsection*{1. Recepción, trazos y preparación}
 \begin{tabularx}{\linewidth}{>{\raggedright\arraybackslash}p{0.93\linewidth} >{\centering\arraybackslash}m{0.07\linewidth}}
 	\toprule
 	\textbf{Actividad a revisar / ordenar} & \textbf{Estado}\\
 	\midrule
 	Viguetas \textbf{íntegras} (sin fisuras/daños); bovedillas completas, secas y limpias. & \cbox \\
 	Trazos de modulación y ejes verificados; \textbf{longitud de apoyo} mínima de viguetas asegurada. & \cbox \\
 	Apoyo perimetral (castillos, dalas, trabes) listo y a nivel. & \cbox \\
 	Frente de trabajo limpio; rutas de circulación y pasarelas definidas. & \cbox \\
 	\bottomrule
 \end{tabularx}
 
 \subsubsection*{2. Montaje de viguetas y bovedillas}
 \begin{tabularx}{\linewidth}{>{\raggedright\arraybackslash}p{0.93\linewidth} >{\centering\arraybackslash}m{0.07\linewidth}}
 	\toprule
 	\textbf{Actividad a revisar / ordenar} & \textbf{Estado}\\
 	\midrule
 	Colocación de viguetas con \textbf{separación de catálogo} y alineación; extremos perfectamente apoyados. & \cbox \\
 	Bovedilla usada como escantillón; \textbf{no forzar} piezas ni dejar huecos sin soporte. & \cbox \\
 	\textbf{Apuntalamiento temporal} continuo y nivelado bajo las viguetas, arriostrado. & \cbox \\
 	Elementos de borde (bezotes/rieles, dalas perimetrales) colocados. & \cbox \\
 	\textit{Orden:} prohibido cortar bovedillas o ranurar viguetas; cualquier hueco debe aprobarse y reforzarse. & \cbox \\
 	\bottomrule
 \end{tabularx}
 
 \subsubsection*{3. Refuerzos y armaduras complementarias}
 \begin{tabularx}{\linewidth}{>{\raggedright\arraybackslash}p{0.93\linewidth} >{\centering\arraybackslash}m{0.07\linewidth}}
 	\toprule
 	\textbf{Actividad a revisar / ordenar} & \textbf{Estado}\\
 	\midrule
 	\textbf{Malla de compresión} (calibre y cuadro) asentada a la \textbf{cota} indicada con separadores. & \cbox \\
 	Refuerzo negativo sobre apoyos y \textbf{perimetrales} (amarres/escuadras) conforme planos o manual. & \cbox \\
 	Refuerzo adicional en zonas de huecos, pasos MEP y cargas puntuales. & \cbox \\
 	Coordinación de instalaciones: tuberías fijas para evitar flotación durante el colado. & \cbox \\
 	\bottomrule
 \end{tabularx}
 
 \subsubsection*{4. Pre-colado}
 \begin{tabularx}{\linewidth}{>{\raggedright\arraybackslash}p{0.93\linewidth} >{\centering\arraybackslash}m{0.07\linewidth}}
 	\toprule
 	\textbf{Actividad a revisar / ordenar} & \textbf{Estado}\\
 	\midrule
 	Humedecer viguetas y bovedillas \textbf{sin encharcar} antes del colado. & \cbox \\
 	Equipo y accesos listos: vibrador (uso \textbf{ligero}), regla/guía, manguera de bomba. & \cbox \\
 	Slump y temperatura del concreto de capa verificados. & \cbox \\
 	Pasarelas para desplazamiento \textbf{sin pisar} directamente las bovedillas. & \cbox \\
 	\bottomrule
 \end{tabularx}
 
 \subsubsection*{5. Colado de la capa de compresión y acabado}
 \begin{tabularx}{\linewidth}{>{\raggedright\arraybackslash}p{0.93\linewidth} >{\centering\arraybackslash}m{0.07\linewidth}}
 	\toprule
 	\textbf{Actividad a revisar / ordenar} & \textbf{Estado}\\
 	\midrule
 	Espesor de capa conforme plano/catálogo; verter por paños y \textbf{vibrado suave} (evitar flotar bovedillas). & \cbox \\
 	Nivelación con regla; acabado especificado (escobillado/liso). & \cbox \\
 	Control de bordes y encuentros con muros/elementos verticales. & \cbox \\
 	Prohibido acopio de materiales sobre la losa fresca; controlar \textbf{cargas de construcción}. & \cbox \\
 	\bottomrule
 \end{tabularx}
 
 \subsubsection*{6. Curado}
 \begin{tabularx}{\linewidth}{>{\raggedright\arraybackslash}p{0.93\linewidth} >{\centering\arraybackslash}m{0.07\linewidth}}
 	\toprule
 	\textbf{Actividad a revisar / ordenar} & \textbf{Estado}\\
 	\midrule
 	Inicio inmediato del curado: \textbf{membrana}, riego o mantas; continuidad durante el periodo indicado. & \cbox \\
 	Proteger del viento/sol/lluvia; restringir tránsito hasta autorización. & \cbox \\
 	\bottomrule
 \end{tabularx}
 
 \subsubsection*{7. Retiro de apuntalamientos y control final}
 \begin{tabularx}{\linewidth}{>{\raggedright\arraybackslash}p{0.93\linewidth} >{\centering\arraybackslash}m{0.07\linewidth}}
 	\toprule
 	\textbf{Actividad a revisar / ordenar} & \textbf{Estado}\\
 	\midrule
 	Retirar \textbf{apuntalamientos} según resistencia alcanzada y \textbf{tablas del fabricante}. & \cbox \\
 	Revisión de superficie: fisuras, flechas, nivelación general; corregir no conformidades. & \cbox \\
 	\bottomrule
 \end{tabularx}
 
 %%%%%%%%%%%%%%%%%%%%%%%%%%%%%%%%%%%%%%%%%%%%%%%%%%%%%%%%%%%%%%%%%%%%
 
 
 
 \section{Losa nervada (aligerada con casetones o bloques de relleno)}
 
 \subsection*{Descripción general}
 La \textbf{losa nervada} es un sistema aligerado de concreto reforzado en el que el peso propio se reduce mediante el uso de casetones o bloques de relleno que generan nervios longitudinales y transversales. Puede vaciarse \emph{in situ} con cimbra modular o moldes reutilizables (plásticos, metálicos o de poliestireno). Combina economía de material con alta rigidez y capacidad para cubrir claros medianos a grandes.
 
 
 
 % =====================================================================
 \subsection{Checklist de supervisión}
 
 \subsubsection*{1. Preparación y montaje de cimbra}
 \begin{tabularx}{\linewidth}{>{\raggedright\arraybackslash}p{0.93\linewidth} >{\centering\arraybackslash}m{0.07\linewidth}}
 	\toprule
 	\textbf{Actividad a revisar / ordenar} & \textbf{Estado}\\
 	\midrule
 	Cimbra nivelada y firme, con bases sólidas y arriostramiento. & \cbox \\
 	Casetones limpios, sin deformaciones y correctamente anclados (sin fugas). & \cbox \\
 	Verificación de contraflecha, si aplica. & \cbox \\
 	Área libre de residuos, con accesos y pasarelas seguras. & \cbox \\
 	\bottomrule
 \end{tabularx}
 
 \subsubsection*{2. Armado de acero}
 \begin{tabularx}{\linewidth}{>{\raggedright\arraybackslash}p{0.93\linewidth} >{\centering\arraybackslash}m{0.07\linewidth}}
 	\toprule
 	\textbf{Actividad a revisar / ordenar} & \textbf{Estado}\\
 	\midrule
 	Barras en nervios principales y secundarios según plano; \textbf{recubrimientos} garantizados. & \cbox \\
 	Refuerzos superiores e inferiores en apoyos, bordes y claros. & \cbox \\
 	Malla de temperatura en capa de compresión, bien apoyada en separadores. & \cbox \\
 	Refuerzos en huecos o pasos de instalaciones. & \cbox \\
 	\bottomrule
 \end{tabularx}
 
 \subsubsection*{3. Colado y vibrado}
 \begin{tabularx}{\linewidth}{>{\raggedright\arraybackslash}p{0.93\linewidth} >{\centering\arraybackslash}m{0.07\linewidth}}
 	\toprule
 	\textbf{Actividad a revisar / ordenar} & \textbf{Estado}\\
 	\midrule
 	Concreto con slump y resistencia especificada; colado por paños uniformes. & \cbox \\
 	Vibrado controlado con aguja fina (evitar flotación de casetones). & \cbox \\
 	Nivelación y acabado de capa de compresión conforme plano. & \cbox \\
 	\bottomrule
 \end{tabularx}
 
 \subsubsection*{4. Curado y retiro de moldes}
 \begin{tabularx}{\linewidth}{>{\raggedright\arraybackslash}p{0.93\linewidth} >{\centering\arraybackslash}m{0.07\linewidth}}
 	\toprule
 	\textbf{Actividad a revisar / ordenar} & \textbf{Estado}\\
 	\midrule
 	Inicio inmediato de curado (membrana, riego o mantas húmedas). & \cbox \\
 	Retiro de casetones cuando el concreto haya alcanzado resistencia mínima. & \cbox \\
 	Revisión visual de superficie y nervios: sin fisuras ni desprendimientos. & \cbox \\
 	\bottomrule
 \end{tabularx}
 
 
 % ================================================================
 \section{Losa colaborante (losacero o steel deck)}
 
 \subsection*{Descripción general}
 La \textbf{losacero} o losa colaborante es un sistema compuesto formado por una lámina metálica acanalada (deck) que actúa como cimbra y, en algunos casos, como refuerzo a tracción una vez endurecido el concreto. Permite una ejecución rápida, sin cimbra tradicional, y brinda resistencia durante el vaciado y al trabajar compuesta con el concreto mediante conectores de cortante.
 
 
 
 % =====================================================================
 \subsection{Checklist de supervisión}
 
 \subsubsection*{1. Preparación y colocación del deck}
 \begin{tabularx}{\linewidth}{>{\raggedright\arraybackslash}p{0.93\linewidth} >{\centering\arraybackslash}m{0.07\linewidth}}
 	\toprule
 	\textbf{Actividad a revisar / ordenar} & \textbf{Estado}\\
 	\midrule
 	Láminas de deck limpias, sin deformaciones ni corrosión; calibre y perfil conforme plano. & \cbox \\
 	\textbf{Fijación} del deck con pernos o tornillos a vigas/canales según diseño. & \cbox \\
 	Solapes y uniones transversales sellados (sin fugas de lechada). & \cbox \\
 	Apoyos perimetrales o ángulos colocados correctamente. & \cbox \\
 	\bottomrule
 \end{tabularx}
 
 \subsubsection*{2. Conectores de cortante y refuerzo}
 \begin{tabularx}{\linewidth}{>{\raggedright\arraybackslash}p{0.93\linewidth} >{\centering\arraybackslash}m{0.07\linewidth}}
 	\toprule
 	\textbf{Actividad a revisar / ordenar} & \textbf{Estado}\\
 	\midrule
 	Conectores de cortante soldados: altura, alineación y fusión visual correctas. & \cbox \\
 	Malla o varillas superiores a recubrimiento correcto desde cresta del deck. & \cbox \\
 	Refuerzo negativo en apoyos y zonas críticas instalado según planos. & \cbox \\
 	\bottomrule
 \end{tabularx}
 
 \subsubsection*{3. Colado y vibrado}
 \begin{tabularx}{\linewidth}{>{\raggedright\arraybackslash}p{0.93\linewidth} >{\centering\arraybackslash}m{0.07\linewidth}}
 	\toprule
 	\textbf{Actividad a revisar / ordenar} & \textbf{Estado}\\
 	\midrule
 	Colado perpendicular a las nervaduras del deck; espesor conforme diseño. & \cbox \\
 	Vibrado superficial moderado (evitar separación de concreto/acero). & \cbox \\
 	Nivelación y acabado uniforme en toda la superficie. & \cbox \\
 	\bottomrule
 \end{tabularx}
 
 \subsubsection*{4. Curado y control final}
 \begin{tabularx}{\linewidth}{>{\raggedright\arraybackslash}p{0.93\linewidth} >{\centering\arraybackslash}m{0.07\linewidth}}
 	\toprule
 	\textbf{Actividad a revisar / ordenar} & \textbf{Estado}\\
 	\midrule
 	Aplicar \textbf{membrana de curado} o riego; proteger de sol y viento. & \cbox \\
 	Verificar espesores y resistencia a compresión del concreto. & \cbox \\
 	Comprobación de uniones, conectores y rigidez del sistema. & \cbox \\
 	\bottomrule
 \end{tabularx}
 
 
 
 
 
 %%aqwuieeeeeeeeeeeeeeeeeeeeeeeeeeeeeeeeeeeeee
 
 
 
 %%%%%%%%%%%%%%%%%%%%%%%%%%%%%%%%%%%%%%%%%%%%%%%%%%%%%%%%%%%%%%%




	%-----------------------------------------------------





%%quieeeeeeee


	
	%--------------SUBTITULOS 2---------------------
     \newpage
	% Segunda introducción
	\newpage
	\stepcounter{intro}
	\section*{\centering \theintro\ Proceso constructivo de elementos de concreto}
	\addcontentsline{toc}{section}{\theintro\ Descripción y proceso constructivo de elementos de concreto}
	\setcounter{section}{0}
	
	
	%REINICIAR CONTEO DE SECTION 
	%\setcounter{section}{0}
	
	
	
	
	%%%%%%%%%%%%%%%::::::::::::::::::::::::::::
	
	
	
	\section{Levantado de muros de mampostería estructural}
	
	\subsection*{Descripción general}
	El \textbf{levantado de muros de mampostería} consiste en la construcción de elementos verticales formados por unidades modulares (bloques de concreto, tabiques cerámicos o piedra), unidos mediante mortero estructural o de pega. Los muros pueden ser portantes o de cerramiento y, en muchos casos, trabajan en conjunto con columnas, castillos y vigas de confinamiento para resistir cargas verticales y laterales.
	

	
	% =====================================================================
	\subsection{Checklist de supervisión}
	
	\subsubsection*{1. Trazado y preparación del área}
	\begin{tabularx}{\linewidth}{>{\raggedright\arraybackslash}p{0.93\linewidth} >{\centering\arraybackslash}m{0.07\linewidth}}
		\toprule
		\textbf{Actividad a revisar / ordenar} & \textbf{Estado}\\
		\midrule
		Ejes y alineamientos del muro trazados conforme planos. & \cbox \\
		Base limpia, nivelada y humedecida antes del asiento del primer bloque. & \cbox \\
		Verificación de espesores y longitudes de muros según planos. & \cbox \\
		Materiales (bloques, mortero, acero) acopiados y protegidos del sol/lluvia. & \cbox \\
		\bottomrule
	\end{tabularx}
	
	\subsubsection*{2. Colocación de bloques y mortero}
	\begin{tabularx}{\linewidth}{>{\raggedright\arraybackslash}p{0.93\linewidth} >{\centering\arraybackslash}m{0.07\linewidth}}
		\toprule
		\textbf{Actividad a revisar / ordenar} & \textbf{Estado}\\
		\midrule
		Primera hilada nivelada y plomeada correctamente. & \cbox \\
		Mortero de pega con consistencia adecuada; juntas horizontales de 1–1.5 cm uniformes. & \cbox \\
		Bloques asentados con junta llena (sin huecos ni faltantes de mortero). & \cbox \\
		Revisión de verticalidad y alineamiento cada tres hiladas. & \cbox \\
		\bottomrule
	\end{tabularx}
	
	\subsubsection*{3. Refuerzo estructural}
	\begin{tabularx}{\linewidth}{>{\raggedright\arraybackslash}p{0.93\linewidth} >{\centering\arraybackslash}m{0.07\linewidth}}
		\toprule
		\textbf{Actividad a revisar / ordenar} & \textbf{Estado}\\
		\midrule
		Refuerzos verticales (barras) colocados en celdas según planos estructurales. & \cbox \\
		Refuerzo horizontal o malla colocado cada cierto número de hiladas (según diseño). & \cbox \\
		Celdas con acero \textbf{rellenas con grout} o mortero fluido hasta la altura indicada. & \cbox \\
		Empalmes de acero respetan longitudes mínimas de traslape. & \cbox \\
		\bottomrule
	\end{tabularx}
	
	\subsubsection*{4. Confinamiento y amarres}
	\begin{tabularx}{\linewidth}{>{\raggedright\arraybackslash}p{0.93\linewidth} >{\centering\arraybackslash}m{0.07\linewidth}}
		\toprule
		\textbf{Actividad a revisar / ordenar} & \textbf{Estado}\\
		\midrule
		Colocación de \textbf{castillos} y \textbf{vigas de amarre} según planos. & \cbox \\
		Amarres entre muros y elementos estructurales garantizados (ganchos, varillas, conectores). & \cbox \\
		Continuidad del refuerzo vertical en castillos con el refuerzo del muro. & \cbox \\
		Conexión muro–columna o muro–viga limpia, sin lechada ni restos sueltos. & \cbox \\
		\bottomrule
	\end{tabularx}
	
	\subsubsection*{5. Acabado y curado}
	\begin{tabularx}{\linewidth}{>{\raggedright\arraybackslash}p{0.93\linewidth} >{\centering\arraybackslash}m{0.07\linewidth}}
		\toprule
		\textbf{Actividad a revisar / ordenar} & \textbf{Estado}\\
		\midrule
		Juntas alineadas y limpias; superficie lista para repello o acabado final. & \cbox \\
		Curado de mortero (mantener humedad durante 3 días como mínimo). & \cbox \\
		Eliminación de mortero sobrante o manchas en superficie. & \cbox \\
		Área segura, libre de materiales sueltos o puntales inestables. & \cbox \\
		\bottomrule
	\end{tabularx}
	
	\subsubsection*{6. Control final}
	\begin{tabularx}{\linewidth}{>{\raggedright\arraybackslash}p{0.93\linewidth} >{\centering\arraybackslash}m{0.07\linewidth}}
		\toprule
		\textbf{Actividad a revisar / ordenar} & \textbf{Estado}\\
		\midrule
		Revisión de verticalidad total y alineamiento del muro terminado. & \cbox \\
		Comprobación de nivel en la parte superior del muro para recibir vigas o losas. & \cbox \\
		Verificación de juntas uniformes, sin grietas visibles ni desviaciones. & \cbox \\
		\bottomrule
	\end{tabularx}
	
	
	
		
	
	%%%%%%%%%%%%%:::::::::::::::::::::::::::::::::::::::::::
	
		
	\section{Construcción de mochetas de concreto}
	
	\subsection*{Descripción general}
	Las \textbf{mochetas} son elementos de concreto armado o de mampostería ubicados a los lados de puertas, ventanas o huecos estructurales, cuya función principal es \textbf{confinar, reforzar y sostener los bordes} del vano. También ayudan a anclar marcos metálicos o de madera y a transmitir las cargas del dintel hacia los muros o columnas. Su correcta ejecución es fundamental para evitar fisuras, asentamientos diferenciales o pérdida de alineación en los vanos.
	

	
	% =====================================================================
	\subsection{Checklist de supervisión}
	
	\subsubsection*{1. Preparación y trazado}
	\begin{tabularx}{\linewidth}{>{\raggedright\arraybackslash}p{0.93\linewidth} >{\centering\arraybackslash}m{0.07\linewidth}}
		\toprule
		\textbf{Actividad a revisar / ordenar} & \textbf{Estado}\\
		\midrule
		Ejes y dimensiones del vano marcados conforme planos. & \cbox \\
		Superficie de contacto del muro limpia y humedecida. & \cbox \\
		Revisión de verticalidad y nivel antes del encofrado. & \cbox \\
		Materiales disponibles (acero, concreto, cimbra) en el frente de trabajo. & \cbox \\
		\bottomrule
	\end{tabularx}
	
	\subsubsection*{2. Armado de refuerzo}
	\begin{tabularx}{\linewidth}{>{\raggedright\arraybackslash}p{0.93\linewidth} >{\centering\arraybackslash}m{0.07\linewidth}}
		\toprule
		\textbf{Actividad a revisar / ordenar} & \textbf{Estado}\\
		\midrule
		Varillas verticales colocadas según plano (diámetro, separación y longitud de anclaje). & \cbox \\
		Estribos o anillos bien amarrados, con recubrimiento mínimo especificado. & \cbox \\
		Empalmes y amarres verificados (sin desplazamiento durante el colado). & \cbox \\
		Refuerzo unido correctamente con el acero del muro o columna adyacente. & \cbox \\
		\bottomrule
	\end{tabularx}
	
	\subsubsection*{3. Encofrado y alineación}
	\begin{tabularx}{\linewidth}{>{\raggedright\arraybackslash}p{0.93\linewidth} >{\centering\arraybackslash}m{0.07\linewidth}}
		\toprule
		\textbf{Actividad a revisar / ordenar} & \textbf{Estado}\\
		\midrule
		Cimbra firme, alineada y estanca; asegurada para evitar fugas de lechada. & \cbox \\
		Plomos y niveles revisados antes del colado. & \cbox \\
		Aplicación uniforme de desmoldante sin contaminar el acero. & \cbox \\
		Pasarelas seguras o plataformas estables en vanos altos. & \cbox \\
		\bottomrule
	\end{tabularx}
	
	\subsubsection*{4. Colado y vibrado}
	\begin{tabularx}{\linewidth}{>{\raggedright\arraybackslash}p{0.93\linewidth} >{\centering\arraybackslash}m{0.07\linewidth}}
		\toprule
		\textbf{Actividad a revisar / ordenar} & \textbf{Estado}\\
		\midrule
		Concreto con slump y resistencia especificada (f’c conforme diseño). & \cbox \\
		Colado continuo y sin interrupciones; vibrado moderado con aguja fina. & \cbox \\
		Buena adherencia entre concreto nuevo y superficie existente. & \cbox \\
		Evitar segregación o desplazamiento del acero durante el vaciado. & \cbox \\
		\bottomrule
	\end{tabularx}
	
	\subsubsection*{5. Curado y desencofrado}
	\begin{tabularx}{\linewidth}{>{\raggedright\arraybackslash}p{0.93\linewidth} >{\centering\arraybackslash}m{0.07\linewidth}}
		\toprule
		\textbf{Actividad a revisar / ordenar} & \textbf{Estado}\\
		\midrule
		Curado iniciado inmediatamente tras el colado (mantas húmedas o riego). & \cbox \\
		Tiempo mínimo de curado cumplido (≥ 7 días o según especificación). & \cbox \\
		Desencofrado sin dañar aristas o caras del elemento. & \cbox \\
		Verificación de plomo, alineación y ancho uniforme después del curado. & \cbox \\
		\bottomrule
	\end{tabularx}
	
	\subsubsection*{6. Acabado y control final}
	\begin{tabularx}{\linewidth}{>{\raggedright\arraybackslash}p{0.93\linewidth} >{\centering\arraybackslash}m{0.07\linewidth}}
		\toprule
		\textbf{Actividad a revisar / ordenar} & \textbf{Estado}\\
		\midrule
		Superficies lisas, sin vacíos ni rebabas, listas para recibir marco o acabado. & \cbox \\
		Dimensiones finales del vano verificadas (altura, ancho y nivel del dintel). & \cbox \\
		Revisión final de alineación y plomo de ambas mochetas. & \cbox \\
		Limpieza del área y retiro de cimbra o residuos. & \cbox \\
		\bottomrule
	\end{tabularx}
	
	
	
	
	
	
	
	
	
	
	%%%%%%%%%%%%%%%%%%%%%%%%%%%%%%%%%%%%%%%%%%%%%%%%%%%
	
	
	
	
	
	\section{Construcción de soleras de concreto armado}
	
	\subsection*{Descripción general}
	Las \textbf{soleras} son elementos horizontales de concreto armado colocados en la parte superior o inferior de los muros, con el propósito de \textbf{amarrar, nivelar y distribuir cargas}.  
	Pueden ser:
	\begin{itemize}[leftmargin=*, itemsep=0pt]
		\item \textbf{Soleras de cimentación o nivelación:} sirven como base nivelada para levantar muros de mampostería.  
		\item \textbf{Soleras de amarre o cerramiento:} se ubican en la parte superior del muro para brindar rigidez, continuidad estructural y anclaje con losas o techos.
	\end{itemize}
	Estas piezas aseguran la integridad del sistema de muros, evitando fisuras por asentamientos diferenciales y mejorando el comportamiento sísmico del conjunto.
	

	
	% =====================================================================
	\subsection{Checklist de supervisión}
	
	\subsubsection*{1. Trazado y preparación del área}
	\begin{tabularx}{\linewidth}{>{\raggedright\arraybackslash}p{0.93\linewidth} >{\centering\arraybackslash}m{0.07\linewidth}}
		\toprule
		\textbf{Actividad a revisar / ordenar} & \textbf{Estado}\\
		\midrule
		Ejes y niveles de solera trazados conforme planos estructurales. & \cbox \\
		Superficie de apoyo limpia, nivelada y humedecida. & \cbox \\
		Verificación del ancho y altura del elemento según diseño. & \cbox \\
		Materiales acopiados (acero, cimbra, concreto, separadores) en buenas condiciones. & \cbox \\
		\bottomrule
	\end{tabularx}
	
	\subsubsection*{2. Colocación del acero de refuerzo}
	\begin{tabularx}{\linewidth}{>{\raggedright\arraybackslash}p{0.93\linewidth} >{\centering\arraybackslash}m{0.07\linewidth}}
		\toprule
		\textbf{Actividad a revisar / ordenar} & \textbf{Estado}\\
		\midrule
		Barras longitudinales colocadas según plano (diámetro, cantidad y recubrimiento). & \cbox \\
		Estribos o anillos instalados con la separación indicada. & \cbox \\
		Anclajes y traslapes de acero correctos en uniones con muros o columnas. & \cbox \\
		Refuerzo libre de grasa, aceite u óxido suelto. & \cbox \\
		\bottomrule
	\end{tabularx}
	
	\subsubsection*{3. Encofrado y nivelación}
	\begin{tabularx}{\linewidth}{>{\raggedright\arraybackslash}p{0.93\linewidth} >{\centering\arraybackslash}m{0.07\linewidth}}
		\toprule
		\textbf{Actividad a revisar / ordenar} & \textbf{Estado}\\
		\midrule
		Cimbra firme, alineada y estanca, con secciones uniformes. & \cbox \\
		Aplicación uniforme de desmoldante (sin contaminar acero). & \cbox \\
		Nivel y plomo de cimbra verificados antes del colado. & \cbox \\
		\bottomrule
	\end{tabularx}
	
	\subsubsection*{4. Colado y vibrado}
	\begin{tabularx}{\linewidth}{>{\raggedright\arraybackslash}p{0.93\linewidth} >{\centering\arraybackslash}m{0.07\linewidth}}
		\toprule
		\textbf{Actividad a revisar / ordenar} & \textbf{Estado}\\
		\midrule
		Concreto con f’c y slump especificados (control de calidad). & \cbox \\
		Colado continuo sin interrupciones, vibrado moderado. & \cbox \\
		Superficie nivelada y sin segregación de materiales. & \cbox \\
		Evitar movimiento del acero o cimbra durante el vaciado. & \cbox \\
		\bottomrule
	\end{tabularx}
	
	\subsubsection*{5. Curado y desencofrado}
	\begin{tabularx}{\linewidth}{>{\raggedright\arraybackslash}p{0.93\linewidth} >{\centering\arraybackslash}m{0.07\linewidth}}
		\toprule
		\textbf{Actividad a revisar / ordenar} & \textbf{Estado}\\
		\midrule
		Iniciar curado inmediatamente tras el acabado (membrana, riego o mantas). & \cbox \\
		Mantener curado continuo por mínimo 7 días o según especificación. & \cbox \\
		Desencofrado sin dañar bordes o caras del elemento. & \cbox \\
		\bottomrule
	\end{tabularx}
	
	\subsubsection*{6. Verificación final}
	\begin{tabularx}{\linewidth}{>{\raggedright\arraybackslash}p{0.93\linewidth} >{\centering\arraybackslash}m{0.07\linewidth}}
		\toprule
		\textbf{Actividad a revisar / ordenar} & \textbf{Estado}\\
		\midrule
		Revisión de nivel y alineación con muros o losas adyacentes. & \cbox \\
		Dimensiones finales verificadas conforme planos. & \cbox \\
		Superficie lista para recibir el muro o el acabado siguiente. & \cbox \\
		\bottomrule
	\end{tabularx}
	
	
	
	
	
	
	
	%%%%%%%%%%%%%%%%%%%%%%%%%%%%%%%%%%%%%%
	
	
	
	\section{Vigas de concreto reforzado}
	
	\subsection*{Descripción general}
	Las \textbf{vigas de concreto reforzado} son elementos lineales que transfieren cargas de losas, muros y otros componentes hacia columnas o apoyos, trabajando principalmente a flexión y cortante. Su desempeño depende de un \emph{encofrado} rígido y alineado, un \emph{armado} conforme a planos (acero longitudinal y estribos), un \emph{colado y vibrado} adecuados y un \emph{curado} oportuno para lograr la resistencia especificada.
	

	% =====================================================================
	\subsection{Checklist de supervisión}
	
	\subsubsection*{1. Trazado, cimbra y apoyos}
	\begin{tabularx}{\linewidth}{>{\raggedright\arraybackslash}p{0.93\linewidth} >{\centering\arraybackslash}m{0.07\linewidth}}
		\toprule
		\textbf{Actividad a revisar / ordenar} & \textbf{Estado}\\
		\midrule
		Ejes, niveles y dimensiones de la viga trazados conforme planos. & \cbox \\
		Cimbra \textbf{alineada, rígida y estanca}, con arriostramiento adecuado. & \cbox \\
		Apoyos y juntas de cimbra sellados (sin fugas de lechada). & \cbox \\
		Contraflecha prevista cuando la especificó el diseñador. & \cbox \\
		Aplicación uniforme de desmoldante (sin tocar el acero). & \cbox \\
		\bottomrule
	\end{tabularx}
	
	\subsubsection*{2. Armado (acero longitudinal y estribos)}
	\begin{tabularx}{\linewidth}{>{\raggedright\arraybackslash}p{0.93\linewidth} >{\centering\arraybackslash}m{0.07\linewidth}}
		\toprule
		\textbf{Actividad a revisar / ordenar} & \textbf{Estado}\\
		\midrule
		Diámetros y número de varillas según plano; \textbf{recubrimientos} garantizados con separadores. & \cbox \\
		\textbf{Estribos} con diámetro y espaciamiento conforme plano (zonas de cortante/ánclaje). & \cbox \\
		Longitudes de \textbf{traslape y anclaje} correctas; ganchos en extremos donde aplique. & \cbox \\
		Amarres firmes del acero para evitar desplazamientos durante el colado. & \cbox \\
		Limpieza del refuerzo (sin grasa ni óxido suelto). & \cbox \\
		\bottomrule
	\end{tabularx}
	
	\subsubsection*{3. Embebidos, huecos y coordinaciones}
	\begin{tabularx}{\linewidth}{>{\raggedright\arraybackslash}p{0.93\linewidth} >{\centering\arraybackslash}m{0.07\linewidth}}
		\toprule
		\textbf{Actividad a revisar / ordenar} & \textbf{Estado}\\
		\midrule
		Mangas, pernos, anclajes y ductos ubicados según plano (sin cortar estribos/longitudinales). & \cbox \\
		Separación mínima entre ductos y acero respetada; sellos para evitar fugas. & \cbox \\
		Coordinación con losas/muros adyacentes (continuidad y anclajes). & \cbox \\
		\bottomrule
	\end{tabularx}
	
	\subsubsection*{4. Pre-colado}
	\begin{tabularx}{\linewidth}{>{\raggedright\arraybackslash}p{0.93\linewidth} >{\centering\arraybackslash}m{0.07\linewidth}}
		\toprule
		\textbf{Actividad a revisar / ordenar} & \textbf{Estado}\\
		\midrule
		Cimbra y acero \textbf{limpios y secos}; retirar residuos y polvo. & \cbox \\
		Vibradores operativos; accesos y mangueras de bomba listos. & \cbox \\
		Verificación de \textbf{slump} y temperatura del concreto inicial. & \cbox \\
		\bottomrule
	\end{tabularx}
	
	\subsubsection*{5. Colado, vibrado y acabado}
	\begin{tabularx}{\linewidth}{>{\raggedright\arraybackslash}p{0.93\linewidth} >{\centering\arraybackslash}m{0.07\linewidth}}
		\toprule
		\textbf{Actividad a revisar / ordenar} & \textbf{Estado}\\
		\midrule
		Colado por capas, evitando caída libre excesiva y segregación. & \cbox \\
		\textbf{Vibrado interno} uniforme: penetración rápida, retiro lento y traslape de zonas. & \cbox \\
		Coronas y caras de viga niveladas; encuentros con losas/muros sin huecos. & \cbox \\
		Control de juntas de construcción (si aplica) en ubicación prevista. & \cbox \\
		\bottomrule
	\end{tabularx}
	
	\subsubsection*{6. Curado y descimbrado}
	\begin{tabularx}{\linewidth}{>{\raggedright\arraybackslash}p{0.93\linewidth} >{\centering\arraybackslash}m{0.07\linewidth}}
		\toprule
		\textbf{Actividad a revisar / ordenar} & \textbf{Estado}\\
		\midrule
		Inicio inmediato del \textbf{curado} (membrana, riego o mantas) tras el acabado. & \cbox \\
		Mantener curado continuo durante el tiempo especificado. & \cbox \\
		\textbf{Descimbrado} solo con autorización y resistencia alcanzada; revisión de aristas. & \cbox \\
		\bottomrule
	\end{tabularx}
	
	\subsubsection*{7. Control final y no conformidades}
	\begin{tabularx}{\linewidth}{>{\raggedright\arraybackslash}p{0.93\linewidth} >{\centering\arraybackslash}m{0.07\linewidth}}
		\toprule
		\textbf{Actividad a revisar / ordenar} & \textbf{Estado}\\
		\midrule
		Inspección de superficie: sin nidos, grietas anómalas ni oquedades. & \cbox \\
		Dimensiones y alineación finales conforme planos; continuidad con elementos adyacentes. & \cbox \\
		Registro fotográfico y corrección de no conformidades detectadas. & \cbox \\
		\bottomrule
	\end{tabularx}
	
	
	%%%%%%%%%%%%%%%%%%%%%%%%%%%%%%%%%%%%%%%%%%%%%%%%%%%%%
	
	
	
	
	\section{Columnas de concreto reforzado}
	
	\subsection*{Descripción general}
	Las \textbf{columnas de concreto reforzado} son elementos verticales que transmiten las cargas de las vigas, losas y muros hacia la cimentación. Trabajan principalmente a compresión, pero también resisten flexión y cortante, por lo que requieren una adecuada colocación del acero longitudinal y transversal (estribos), un buen alineamiento y un colado controlado. Son componentes esenciales para la estabilidad y el comportamiento sísmico de la estructura.
	

	
	% =====================================================================
	\subsection{Checklist de supervisión}
	
	\subsubsection*{1. Trazado y preparación}
	\begin{tabularx}{\linewidth}{>{\raggedright\arraybackslash}p{0.93\linewidth} >{\centering\arraybackslash}m{0.07\linewidth}}
		\toprule
		\textbf{Actividad a revisar / ordenar} & \textbf{Estado}\\
		\midrule
		Ubicación exacta de ejes y dimensiones de columnas conforme planos. & \cbox \\
		Nivel y aplomo de plantilla o arranque verificados. & \cbox \\
		Superficie limpia y humedecida antes de colocar acero. & \cbox \\
		\bottomrule
	\end{tabularx}
	
	\subsubsection*{2. Armado de acero}
	\begin{tabularx}{\linewidth}{>{\raggedright\arraybackslash}p{0.93\linewidth} >{\centering\arraybackslash}m{0.07\linewidth}}
		\toprule
		\textbf{Actividad a revisar / ordenar} & \textbf{Estado}\\
		\midrule
		Cantidad, diámetro y separación de barras conforme planos estructurales. & \cbox \\
		Estribos con espaciamiento y diámetro especificado; ganchos cerrados a 135°. & \cbox \\
		Empalmes y traslapes ejecutados con la longitud requerida. & \cbox \\
		Refuerzo limpio, alineado y asegurado para evitar desplazamientos. & \cbox \\
		Separadores colocados para garantizar recubrimiento uniforme. & \cbox \\
		\bottomrule
	\end{tabularx}
	
	\subsubsection*{3. Encofrado y alineación}
	\begin{tabularx}{\linewidth}{>{\raggedright\arraybackslash}p{0.93\linewidth} >{\centering\arraybackslash}m{0.07\linewidth}}
		\toprule
		\textbf{Actividad a revisar / ordenar} & \textbf{Estado}\\
		\midrule
		Cimbra firme, alineada y arriostrada; sin fugas en juntas. & \cbox \\
		Plomos y escuadras verificados antes del colado. & \cbox \\
		Aplicación uniforme de desmoldante sin contaminar el acero. & \cbox \\
		Ventanas o mirillas de inspección habilitadas si la columna es alta. & \cbox \\
		\bottomrule
	\end{tabularx}
	
	\subsubsection*{4. Colado y vibrado}
	\begin{tabularx}{\linewidth}{>{\raggedright\arraybackslash}p{0.93\linewidth} >{\centering\arraybackslash}m{0.07\linewidth}}
		\toprule
		\textbf{Actividad a revisar / ordenar} & \textbf{Estado}\\
		\midrule
		Concreto con slump y f’c especificado; verificación de temperatura inicial. & \cbox \\
		Colado en capas de 30–50 cm; vibrado interno con aguja fina (sin sobre-vibrar). & \cbox \\
		Evitar segregación y desplazamiento del acero durante el vaciado. & \cbox \\
		Nivel superior preparado para unión con viga o losa. & \cbox \\
		\bottomrule
	\end{tabularx}
	
	\subsubsection*{5. Curado y descimbrado}
	\begin{tabularx}{\linewidth}{>{\raggedright\arraybackslash}p{0.93\linewidth} >{\centering\arraybackslash}m{0.07\linewidth}}
		\toprule
		\textbf{Actividad a revisar / ordenar} & \textbf{Estado}\\
		\midrule
		Iniciar curado inmediato (membrana, riego o mantas húmedas). & \cbox \\
		Mantener humedad continua por mínimo 7 días o según especificación. & \cbox \\
		Desencofrado sin golpear ni dañar las aristas. & \cbox \\
		\bottomrule
	\end{tabularx}
	
	\subsubsection*{6. Control final}
	\begin{tabularx}{\linewidth}{>{\raggedright\arraybackslash}p{0.93\linewidth} >{\centering\arraybackslash}m{0.07\linewidth}}
		\toprule
		\textbf{Actividad a revisar / ordenar} & \textbf{Estado}\\
		\midrule
		Verificación de plomo, alineación y sección transversal. & \cbox \\
		Superficie sin nidos, fisuras ni oquedades visibles. & \cbox \\
		Dimensiones finales y nivel de coronación conforme plano. & \cbox \\
		Limpieza del área y retiro de residuos o cimbra sobrante. & \cbox \\
		\bottomrule
	\end{tabularx}
	
	
	
	%%%%%%%%%%%%%%%%%%%%%%%%%%%%%%%%%%%%%%%%%%%%%%%%%%%%%%%%%
	
	
	\section{Zapata aislada de concreto reforzado}
	
	\subsection*{Descripción general}
	La \textbf{zapata aislada} es un elemento de cimentación superficial de concreto reforzado que transmite la carga de una columna o pilar al terreno. Su desempeño depende de: correcta \emph{excavación y preparación de base}, \emph{armado y recubrimientos}, \emph{colado y vibrado} adecuados, y \emph{curado} oportuno, además de un \emph{relleno y compactación} controlados alrededor de la cimentación.
	

	% =====================================================================
	\subsection{Checklist de supervisión}
	
	\subsubsection*{1. Replanteo, excavación y base}
	\begin{tabularx}{\linewidth}{>{\raggedright\arraybackslash}p{0.93\linewidth} >{\centering\arraybackslash}m{0.07\linewidth}}
		\toprule
		\textbf{Actividad a revisar / ordenar} & \textbf{Estado}\\
		\midrule
		Ejes y dimensiones de zapata trazados conforme planos (ubicación de columna). & \cbox \\
		Cota de desplante alcanzada; estrato competente verificado (sin material suelto). & \cbox \\
		Fondo de excavación nivelado y compactado; sin agua acumulada ni lodos. & \cbox \\
		Taludes o entibaciones seguras y estables (si aplica). & \cbox \\
		\bottomrule
	\end{tabularx}
	
	\subsubsection*{2. Plantilla de limpieza y preparación}
	\begin{tabularx}{\linewidth}{>{\raggedright\arraybackslash}p{0.93\linewidth} >{\centering\arraybackslash}m{0.07\linewidth}}
		\toprule
		\textbf{Actividad a revisar / ordenar} & \textbf{Estado}\\
		\midrule
		Colocación de \textbf{plantilla} de concreto pobre o mortero (espesor y nivel). & \cbox \\
		Superficie limpia, sin polvo ni contaminantes antes del armado. & \cbox \\
		Referencias de nivel y ejes controladas para centrar el armado. & \cbox \\
		\bottomrule
	\end{tabularx}
	
	\subsubsection*{3. Armado y recubrimientos}
	\begin{tabularx}{\linewidth}{>{\raggedright\arraybackslash}p{0.93\linewidth} >{\centering\arraybackslash}m{0.07\linewidth}}
		\toprule
		\textbf{Actividad a revisar / ordenar} & \textbf{Estado}\\
		\midrule
		Armadura inferior en dos direcciones (diámetro, separación y desarrollo) según plano. & \cbox \\
		\textbf{Separadores} suficientes para recubrimiento mínimo en fondo y laterales. & \cbox \\
		Pernos/plantillas de anclaje de columna \textbf{alineados y nivelados}. & \cbox \\
		Empalmes con longitudes de traslape correctas; acero limpio (sin grasa/óxido suelto). & \cbox \\
		\bottomrule
	\end{tabularx}
	
	\subsubsection*{4. Encofrado (si aplica) y control previo al colado}
	\begin{tabularx}{\linewidth}{>{\raggedright\arraybackslash}p{0.93\linewidth} >{\centering\arraybackslash}m{0.07\linewidth}}
		\toprule
		\textbf{Actividad a revisar / ordenar} & \textbf{Estado}\\
		\midrule
		Cimbra para cantos \textbf{alineada y estanca}; juntas selladas (sin fugas). & \cbox \\
		Desmoldante aplicado correctamente (sin contaminar el acero). & \cbox \\
		Verificación de \textbf{slump} y temperatura del concreto de la primera olla. & \cbox \\
		Accesos, vibradores y herramientas listos; ruta de descarga definida. & \cbox \\
		\bottomrule
	\end{tabularx}
	
	\subsubsection*{5. Colado, vibrado y acabado}
	\begin{tabularx}{\linewidth}{>{\raggedright\arraybackslash}p{0.93\linewidth} >{\centering\arraybackslash}m{0.07\linewidth}}
		\toprule
		\textbf{Actividad a revisar / ordenar} & \textbf{Estado}\\
		\midrule
		Concreto con f’c especificado; colado por capas evitando caída libre excesiva. & \cbox \\
		\textbf{Vibrado interno} uniforme (sin sobre-vibrar ni segregar); relleno de esquinas. & \cbox \\
		Nivelación de la \textbf{corona} para recibir columna/arranque; juntas de construcción planificadas. & \cbox \\
		Limpieza de pernos/anclajes; protección contra desplazamientos durante el colado. & \cbox \\
		\bottomrule
	\end{tabularx}
	
	\subsubsection*{6. Curado y desencofrado}
	\begin{tabularx}{\linewidth}{>{\raggedright\arraybackslash}p{0.93\linewidth} >{\centering\arraybackslash}m{0.07\linewidth}}
		\toprule
		\textbf{Actividad a revisar / ordenar} & \textbf{Estado}\\
		\midrule
		Inicio inmediato del \textbf{curado} (membrana, riego o mantas húmedas). & \cbox \\
		Mantener curado continuo por el tiempo especificado (típ. \(\geq\) 7 días). & \cbox \\
		Retiro de cimbra (si existe) sin dañar aristas ni la superficie. & \cbox \\
		\bottomrule
	\end{tabularx}
	
	\subsubsection*{7. Relleno y compactación}
	\begin{tabularx}{\linewidth}{>{\raggedright\arraybackslash}p{0.93\linewidth} >{\centering\arraybackslash}m{0.07\linewidth}}
		\toprule
		\textbf{Actividad a revisar / ordenar} & \textbf{Estado}\\
		\midrule
		Relleno por \textbf{capas delgadas} (\(\leq\) 20--25\,cm), con compactación por capa. & \cbox \\
		Control de \textbf{densidad} (ensayes in situ si el proyecto lo exige). & \cbox \\
		Drenaje perimetral o subdrenes operativos (si fueron especificados). & \cbox \\
		\bottomrule
	\end{tabularx}
	
	\subsubsection*{8. Control final}
	\begin{tabularx}{\linewidth}{>{\raggedright\arraybackslash}p{0.93\linewidth} >{\centering\arraybackslash}m{0.07\linewidth}}
		\toprule
		\textbf{Actividad a revisar / ordenar} & \textbf{Estado}\\
		\midrule
		Dimensiones finales (largo, ancho, espesor) y \textbf{cotas} conforme planos. & \cbox \\
		Posición y verticalidad del arranque/columna respecto a ejes. & \cbox \\
		Superficie sin nidos, oquedades ni segregación visible. & \cbox \\
		Limpieza del área y retiro de excedentes o escombros. & \cbox \\
		\bottomrule
	\end{tabularx}
	
	
	%%%%%%%%%%%%%%%%%%%%%%%%%%%%%%%%%%%%%%%%%%%%%%%%%%%%%%%%%
	
	
	
	
	\section{Ménsulas de concreto reforzado}
	
	\subsection*{Descripción general}
	Las \textbf{ménsulas} o \textbf{voladizos cortos} son elementos estructurales en voladizo de concreto reforzado que se proyectan desde una columna, muro o viga principal para sostener cargas concentradas, tales como vigas secundarias, balcones, losas, cornisas o apoyos metálicos.  
	Su diseño está dominado por esfuerzos de flexión y cortante, por lo que requieren un adecuado anclaje del acero dentro del elemento principal, así como una correcta ejecución de cimbra, armado, colado y curado.
	

	
	% =====================================================================
	\subsection{Checklist de supervisión}
	
	\subsubsection*{1. Replanteo y preparación del área}
	\begin{tabularx}{\linewidth}{>{\raggedright\arraybackslash}p{0.93\linewidth} >{\centering\arraybackslash}m{0.07\linewidth}}
		\toprule
		\textbf{Actividad a revisar / ordenar} & \textbf{Estado}\\
		\midrule
		Ejes, dimensiones y ubicación de la ménsula conforme planos. & \cbox \\
		Nivel de arranque verificado respecto al elemento principal. & \cbox \\
		Superficie de unión (columna o muro) limpia, rugosa y humedecida. & \cbox \\
		Revisión de interferencias con ductos, pernos o embebidos cercanos. & \cbox \\
		\bottomrule
	\end{tabularx}
	
	\subsubsection*{2. Armado de refuerzo}
	\begin{tabularx}{\linewidth}{>{\raggedright\arraybackslash}p{0.93\linewidth} >{\centering\arraybackslash}m{0.07\linewidth}}
		\toprule
		\textbf{Actividad a revisar / ordenar} & \textbf{Estado}\\
		\midrule
		Barras principales colocadas según plano (diámetro, cantidad y longitud de anclaje). & \cbox \\
		Refuerzo de anclaje correctamente embebido dentro de columna/viga. & \cbox \\
		Estribos o refuerzo transversal espaciados conforme plano (zonas de cortante). & \cbox \\
		Separadores colocados para garantizar recubrimientos mínimos. & \cbox \\
		Refuerzo limpio y bien amarrado (sin grasa, pintura u óxido suelto). & \cbox \\
		\bottomrule
	\end{tabularx}
	
	\subsubsection*{3. Encofrado y soportes}
	\begin{tabularx}{\linewidth}{>{\raggedright\arraybackslash}p{0.93\linewidth} >{\centering\arraybackslash}m{0.07\linewidth}}
		\toprule
		\textbf{Actividad a revisar / ordenar} & \textbf{Estado}\\
		\midrule
		Cimbra firme, nivelada y estanca, con refuerzo inferior temporal. & \cbox \\
		Arriostramientos o soportes capaces de resistir el peso del colado. & \cbox \\
		Aplicación uniforme de desmoldante (sin tocar el acero). & \cbox \\
		Alineación y nivelación de la cara superior y bordes. & \cbox \\
		\bottomrule
	\end{tabularx}
	
	\subsubsection*{4. Colado y vibrado}
	\begin{tabularx}{\linewidth}{>{\raggedright\arraybackslash}p{0.93\linewidth} >{\centering\arraybackslash}m{0.07\linewidth}}
		\toprule
		\textbf{Actividad a revisar / ordenar} & \textbf{Estado}\\
		\midrule
		Concreto con f’c y slump especificado (según diseño). & \cbox \\
		Colado en capas cortas, vibrado con aguja fina sin segregar. & \cbox \\
		Evitar desplazamiento del acero principal durante el vibrado. & \cbox \\
		Nivelación de la superficie superior y acabado adecuado. & \cbox \\
		\bottomrule
	\end{tabularx}
	
	\subsubsection*{5. Curado y retiro de cimbra}
	\begin{tabularx}{\linewidth}{>{\raggedright\arraybackslash}p{0.93\linewidth} >{\centering\arraybackslash}m{0.07\linewidth}}
		\toprule
		\textbf{Actividad a revisar / ordenar} & \textbf{Estado}\\
		\midrule
		Curado inmediato mediante riego o aplicación de membrana. & \cbox \\
		Mantener curado constante durante mínimo 7 días. & \cbox \\
		Retiro de cimbra sólo tras alcanzar resistencia suficiente. & \cbox \\
		Inspección visual de aristas y caras después del desencofrado. & \cbox \\
		\bottomrule
	\end{tabularx}
	
	\subsubsection*{6. Control final}
	\begin{tabularx}{\linewidth}{>{\raggedright\arraybackslash}p{0.93\linewidth} >{\centering\arraybackslash}m{0.07\linewidth}}
		\toprule
		\textbf{Actividad a revisar / ordenar} & \textbf{Estado}\\
		\midrule
		Verificación de dimensiones, alineación y nivel final. & \cbox \\
		Comprobación del anclaje efectivo con el elemento principal. & \cbox \\
		Superficie sin nidos, fisuras ni segregación visible. & \cbox \\
		Área limpia y lista para recibir cargas o acabados. & \cbox \\
		\bottomrule
	\end{tabularx}
	
	
	
	
	
	
	
	
	%%%%%%%%%%%%%%%%%%%%%%%%%%%%%%%%%%%%%%%%%%%%%%%%%%%%%
	
	  \newpage
	% Segunda introducción
	\newpage
	\stepcounter{intro}
	\section*{\centering \theintro\ Descripción y proceso constructivo de drenajes}
	\addcontentsline{toc}{section}{\theintro\ Descripción y proceso constructivo de elementos de concreto}
	\setcounter{section}{0}
	
	
	
	
	\section{Instalación de tuberías (hidrosanitarias)}
	
	\subsection*{Descripción general}
	La \textbf{instalación de tuberías} comprende el trazado, tendido, anclaje, pruebas y protección de redes hidráulicas, sanitarias, pluviales o eléctricas (conduit), ya sean \textbf{embebidas en losa/muro} o \textbf{enterradas en zanja}. La calidad depende de una buena planificación (planos, pendientes, diámetros), correcta sujeción/soporte, sellos estancos en pasos y \textbf{pruebas de presión/estanqueidad} antes del cierre o relleno.
	
	
	
	% =====================================================================
	\subsection{Checklist de supervisión}
	
	\subsubsection*{1. Planificación, trazos y coordinación}
	\begin{tabularx}{\linewidth}{>{\raggedright\arraybackslash}p{0.93\linewidth} >{\centering\arraybackslash}m{0.07\linewidth}}
		\toprule
		\textbf{Actividad a revisar / ordenar} & \textbf{Estado}\\
		\midrule
		Planos actualizados y \textbf{rutas} definidas (diámetros, pendientes, cotas y válvulas). & \cbox \\
		Coordinación con estructura (vigas, losas, muros) y otras instalaciones (eléctrica/HVAC). & \cbox \\
		Trazos sobre obra: ejes, pendientes y niveles marcados y verificados. & \cbox \\
		Autorización para \textbf{perforaciones/pasos} y uso de mangas o casquillos. & \cbox \\
		\bottomrule
	\end{tabularx}
	
	\subsubsection*{2. Materiales, almacenamiento y preparación}
	\begin{tabularx}{\linewidth}{>{\raggedright\arraybackslash}p{0.93\linewidth} >{\centering\arraybackslash}m{0.07\linewidth}}
		\toprule
		\textbf{Actividad a revisar / ordenar} & \textbf{Estado}\\
		\midrule
		Tuberías y accesorios del \textbf{material especificado} (PVC/CPVC/PEAD/galv./conduit) sin daños. & \cbox \\
		Adhesivos, soldaduras y empaques dentro de \textbf{fecha y ficha técnica}. & \cbox \\
		Protección de materiales contra sol/impactos; cortes limpios y biselados donde aplique. & \cbox \\
		\bottomrule
	\end{tabularx}
	
	\subsubsection*{3. Instalaciones \textit{embebidas} en losa/muro (antes del colado)}
	\begin{tabularx}{\linewidth}{>{\raggedright\arraybackslash}p{0.93\linewidth} >{\centering\arraybackslash}m{0.07\linewidth}}
		\toprule
		\textbf{Actividad a revisar / ordenar} & \textbf{Estado}\\
		\midrule
		Tuberías \textbf{sujetas} a armaduras o cimbra para evitar flotación durante el colado. & \cbox \\
		Respeto de \textbf{recubrimientos} y no interferencia con acero de refuerzo. & \cbox \\
		Colocación de cajas, derivaciones y \textbf{mangas} en la posición y cota indicadas. & \cbox \\
		Sellos de pasos y \textbf{sistemas corta–fuego} (si aplica) colocados. & \cbox \\
		\bottomrule
	\end{tabularx}
	
	\subsubsection*{4. Instalaciones \textit{enterradas} en zanja}
	\begin{tabularx}{\linewidth}{>{\raggedright\arraybackslash}p{0.93\linewidth} >{\centering\arraybackslash}m{0.07\linewidth}}
		\toprule
		\textbf{Actividad a revisar / ordenar} & \textbf{Estado}\\
		\midrule
		Profundidad, \textbf{pendiente} y ancho de zanja conforme plano y especificación. & \cbox \\
		\textbf{Cama de apoyo} (arena/grava fina) nivelada y libre de puntas o piedras. & \cbox \\
		Alineación y \textbf{separación} mínima de otras redes; señalización y seguridad de zanja. & \cbox \\
		Relleno por capas y \textbf{compactación} (densidad especificada). & \cbox \\
		\bottomrule
	\end{tabularx}
	
	\subsubsection*{5. Montaje, uniones y soportes (toda condición)}
	\begin{tabularx}{\linewidth}{>{\raggedright\arraybackslash}p{0.93\linewidth} >{\centering\arraybackslash}m{0.07\linewidth}}
		\toprule
		\textbf{Actividad a revisar / ordenar} & \textbf{Estado}\\
		\midrule
		Uniones según sistema: \textbf{cementado}, roscado, termofusión, bridas o press–fit. & \cbox \\
		\textbf{Abrazaderas/soportes} a espaciamientos de catálogo; anclajes en cambios de dirección. & \cbox \\
		Aislamiento térmico/acústico donde aplique; protección UV para plásticos expuestos. & \cbox \\
		Rotulación de redes (agua fría/caliente, sanitaria, pluvial, eléctrica). & \cbox \\
		\bottomrule
	\end{tabularx}
	
	\subsubsection*{6. Pruebas y verificaciones antes del cierre}
	\begin{tabularx}{\linewidth}{>{\raggedright\arraybackslash}p{0.93\linewidth} >{\centering\arraybackslash}m{0.07\linewidth}}
		\toprule
		\textbf{Actividad a revisar / ordenar} & \textbf{Estado}\\
		\midrule
		\textbf{Prueba hidrostática} (hidráulica): presión/tiempo según especificación, sin fugas. & \cbox \\
		\textbf{Prueba de estanqueidad} (sanitaria/pluvial) o continui\-dad en conduits eléctricos. & \cbox \\
		Revisión de pendientes con nivel/laser; purgas y válvulas operativas. & \cbox \\
		Registro fotográfico de rutas y profundidades antes de tapar. & \cbox \\
		\bottomrule
	\end{tabularx}
	
	\subsubsection*{7. Cierre, protección y entrega}
	\begin{tabularx}{\linewidth}{>{\raggedright\arraybackslash}p{0.93\linewidth} >{\centering\arraybackslash}m{0.07\linewidth}}
		\toprule
		\textbf{Actividad a revisar / ordenar} & \textbf{Estado}\\
		\midrule
		Cierre de rozas con mortero; para losa/muro, colado sin desplazar tuberías. & \cbox \\
		Etiquetado y planos \textbf{as-built} actualizados. & \cbox \\
		Áreas limpias y libres de obstrucciones; protección temporal donde transite personal. & \cbox \\
		\bottomrule
	\end{tabularx}
	
	
	
	%%%%%%%%%%%%%%%%%%%%%%%%%%%%%%%%%%%%%%%%%%%%%%%%%%%%%%
	
	
	
	\section{Construcción de pozos de visita}
	
	\subsection*{Descripción general}
	Los \textbf{pozos de visita} son estructuras de concreto (prefabricadas o coladas en sitio) que permiten el acceso, inspección, limpieza y mantenimiento de las redes subterráneas de alcantarillado sanitario, pluvial o drenaje.  
	Generalmente se construyen con una base o fondo (solera), muros cilíndricos o rectangulares y una tapa o losa superior con marco y brocal metálico.  
	Su correcta ejecución garantiza la hermeticidad, alineación con la red y seguridad del personal de mantenimiento.
	

	% =====================================================================
	\subsection{Checklist de supervisión}
	
	\subsubsection*{1. Replanteo, excavación y seguridad}
	\begin{tabularx}{\linewidth}{>{\raggedright\arraybackslash}p{0.93\linewidth} >{\centering\arraybackslash}m{0.07\linewidth}}
		\toprule
		\textbf{Actividad a revisar / ordenar} & \textbf{Estado}\\
		\midrule
		Ubicación y profundidad del pozo conforme planos y cotas de red. & \cbox \\
		Taludes estables o entibación instalada correctamente (seguridad del personal). & \cbox \\
		Base de excavación limpia, nivelada y libre de agua o material suelto. & \cbox \\
		Control de acceso, señalización y ventilación en área de trabajo. & \cbox \\
		\bottomrule
	\end{tabularx}
	
	\subsubsection*{2. Base o solera del pozo}
	\begin{tabularx}{\linewidth}{>{\raggedright\arraybackslash}p{0.93\linewidth} >{\centering\arraybackslash}m{0.07\linewidth}}
		\toprule
		\textbf{Actividad a revisar / ordenar} & \textbf{Estado}\\
		\midrule
		Colocación de capa base o plantilla de concreto pobre nivelada. & \cbox \\
		Solera de concreto reforzado con pendiente hacia el centro. & \cbox \\
		Curado adecuado antes de continuar con los muros. & \cbox \\
		\bottomrule
	\end{tabularx}
	
	\subsubsection*{3. Construcción de muros del pozo}
	\begin{tabularx}{\linewidth}{>{\raggedright\arraybackslash}p{0.93\linewidth} >{\centering\arraybackslash}m{0.07\linewidth}}
		\toprule
		\textbf{Actividad a revisar / ordenar} & \textbf{Estado}\\
		\midrule
		Material de muro conforme especificación (bloque, anillo o concreto vaciado). & \cbox \\
		Verticalidad, espesor y diámetro del pozo verificados. & \cbox \\
		Juntas selladas con mortero impermeable o sello flexible. & \cbox \\
		Instalación de \textbf{escalones metálicos o peldaños} empotrados a distancias iguales. & \cbox \\
		\bottomrule
	\end{tabularx}
	
	\subsubsection*{4. Conexión de tuberías}
	\begin{tabularx}{\linewidth}{>{\raggedright\arraybackslash}p{0.93\linewidth} >{\centering\arraybackslash}m{0.07\linewidth}}
		\toprule
		\textbf{Actividad a revisar / ordenar} & \textbf{Estado}\\
		\midrule
		Entradas y salidas empotradas conforme cotas de diseño y pendiente del colector. & \cbox \\
		Sellado perimetral con mortero o junta de caucho flexible. & \cbox \\
		Canaleta interior (“cuneta”) con pendiente suave, pulida y sin discontinuidades. & \cbox \\
		Verificación de flujo continuo y sin obstrucciones. & \cbox \\
		\bottomrule
	\end{tabularx}
	
	\subsubsection*{5. Losa de tapa y brocal}
	\begin{tabularx}{\linewidth}{>{\raggedright\arraybackslash}p{0.93\linewidth} >{\centering\arraybackslash}m{0.07\linewidth}}
		\toprule
		\textbf{Actividad a revisar / ordenar} & \textbf{Estado}\\
		\midrule
		Losa de tapa colada con refuerzo y espesor especificado. & \cbox \\
		Marco y tapa de fundición alineados con nivel de rasante o pavimento. & \cbox \\
		Orificio central con diámetro de acceso suficiente y acabado limpio. & \cbox \\
		Anclajes y nivelación del brocal asegurados antes del fraguado. & \cbox \\
		\bottomrule
	\end{tabularx}
	
	\subsubsection*{6. Curado, acabado y limpieza final}
	\begin{tabularx}{\linewidth}{>{\raggedright\arraybackslash}p{0.93\linewidth} >{\centering\arraybackslash}m{0.07\linewidth}}
		\toprule
		\textbf{Actividad a revisar / ordenar} & \textbf{Estado}\\
		\midrule
		Curado del concreto de base, muros y tapa durante al menos 7 días. & \cbox \\
		Eliminación de rebabas, residuos o mortero suelto en interior del pozo. & \cbox \\
		Aplicación de sellador o recubrimiento impermeabilizante (si está especificado). & \cbox \\
		Limpieza general y verificación de acceso libre antes de la entrega. & \cbox \\
		\bottomrule
	\end{tabularx}
	
	\subsubsection*{7. Control final y pruebas}
	\begin{tabularx}{\linewidth}{>{\raggedright\arraybackslash}p{0.93\linewidth} >{\centering\arraybackslash}m{0.07\linewidth}}
		\toprule
		\textbf{Actividad a revisar / ordenar} & \textbf{Estado}\\
		\midrule
		Prueba de estanqueidad (llenado con agua) o inspección visual de fugas. & \cbox \\
		Comprobación de cotas de entrada y salida según plano. & \cbox \\
		Revisión de alineación con la red y continuidad del flujo. & \cbox \\
		Documentación y registro fotográfico del pozo terminado. & \cbox \\
		\bottomrule
	\end{tabularx}
	
	
	%%%%%%%%%%%%%%%%%%%%%%%%%%%%%%%%%%%%%%%%%%%%%%%%%%
	
	
	\section{Construcción de pozos de absorción}
	
	\subsection*{Descripción general}
	Los \textbf{pozos de absorción} (o pozos de infiltración) son estructuras verticales que \textbf{disipan y recargan} al subsuelo los caudales de agua (típicamente pluvial o excedentes tratados), mediante paredes y fondo permeables. Su desempeño depende de la \textbf{permeabilidad del estrato}, adecuada \textbf{profundidad y área de infiltración}, y de un \textbf{pretratamiento} que limite el ingreso de sedimentos y grasas.
	
	
	% =====================================================================
	\subsection{Checklist de supervisión}
	
	\subsubsection*{1. Estudios previos y trazos}
	\begin{tabularx}{\linewidth}{>{\raggedright\arraybackslash}p{0.93\linewidth} >{\centering\arraybackslash}m{0.07\linewidth}}
		\toprule
		\textbf{Actividad a revisar / ordenar} & \textbf{Estado}\\
		\midrule
		Resultados de \textbf{prueba de infiltración} y verificación de estratos (permeables en la profundidad de diseño). & \cbox \\
		Ubicación conforme planos: distancias mínimas a cimentaciones, linderos y pozos de agua. & \cbox \\
		Trazos de diámetro y cota de fondo definidos; accesos y zona de acopio seguros. & \cbox \\
		\bottomrule
	\end{tabularx}
	
	\subsubsection*{2. Excavación y seguridad}
	\begin{tabularx}{\linewidth}{>{\raggedright\arraybackslash}p{0.93\linewidth} >{\centering\arraybackslash}m{0.07\linewidth}}
		\toprule
		\textbf{Actividad a revisar / ordenar} & \textbf{Estado}\\
		\midrule
		Excavación vertical a diámetro/cota de proyecto, con \textbf{entibación} o taludes estables. & \cbox \\
		Fondo limpio, sin finos sueltos ni agua encharcada; bombeo (si necesario) operativo. & \cbox \\
		Control de acceso y ventilación del área; señalización y EPP del personal. & \cbox \\
		\bottomrule
	\end{tabularx}
	
	\subsubsection*{3. Lecho drenante y geotextil}
	\begin{tabularx}{\linewidth}{>{\raggedright\arraybackslash}p{0.93\linewidth} >{\centering\arraybackslash}m{0.07\linewidth}}
		\toprule
		\textbf{Actividad a revisar / ordenar} & \textbf{Estado}\\
		\midrule
		Colocación de \textbf{grava limpia} (espesor y granulometría especificados) en el fondo. & \cbox \\
		Instalación de \textbf{geotextil} separador (si aplica) sobre el lecho y/o perimetral. & \cbox \\
		Nivelación del lecho para recibir el cuerpo del pozo. & \cbox \\
		\bottomrule
	\end{tabularx}
	
	\subsubsection*{4. Cuerpo del pozo (elementos permeables)}
	\begin{tabularx}{\linewidth}{>{\raggedright\arraybackslash}p{0.93\linewidth} >{\centering\arraybackslash}m{0.07\linewidth}}
		\toprule
		\textbf{Actividad a revisar / ordenar} & \textbf{Estado}\\
		\midrule
		Anillos perforados / block calado / tubo ranurado \textbf{alineados y plomados}. & \cbox \\
		Perforaciones o ranuras \textbf{libres de obstrucción}; continuidad vertical asegurada. & \cbox \\
		\textbf{Relleno anular} con grava uniforme; envoltura con geotextil si fue especificado. & \cbox \\
		\bottomrule
	\end{tabularx}
	
	\subsubsection*{5. Conexiones, pretratamiento y ventilación}
	\begin{tabularx}{\linewidth}{>{\raggedright\arraybackslash}p{0.93\linewidth} >{\centering\arraybackslash}m{0.07\linewidth}}
		\toprule
		\textbf{Actividad a revisar / ordenar} & \textbf{Estado}\\
		\midrule
		\textbf{Caja trampa de sedimentos/grasas} o filtro previo instalado y accesible. & \cbox \\
		Conexión de tubería de entrada \textbf{con caída} adecuada y sello perimetral estanco. & \cbox \\
		\textbf{Ventilación} del pozo (tubo o respiradero) habilitada y protegida. & \cbox \\
		\bottomrule
	\end{tabularx}
	
	\subsubsection*{6. Tapa, brocal y acceso de mantenimiento}
	\begin{tabularx}{\linewidth}{>{\raggedright\arraybackslash}p{0.93\linewidth} >{\centering\arraybackslash}m{0.07\linewidth}}
		\toprule
		\textbf{Actividad a revisar / ordenar} & \textbf{Estado}\\
		\midrule
		Ejecución de losa/brocal con \textbf{registro} (diámetro suficiente) y marco estable. & \cbox \\
		Cubierta \textbf{cerrada y segura} (evitar ingreso de sólidos y riesgos de caída). & \cbox \\
		Acabado perimetral y pendiente exterior alejando escorrentías superficiales. & \cbox \\
		\bottomrule
	\end{tabularx}
	
	\subsubsection*{7. Puesta en servicio y control final}
	\begin{tabularx}{\linewidth}{>{\raggedright\arraybackslash}p{0.93\linewidth} >{\centering\arraybackslash}m{0.07\linewidth}}
		\toprule
		\textbf{Actividad a revisar / ordenar} & \textbf{Estado}\\
		\midrule
		\textbf{Prueba de infiltración} final (llenado/medición de abatimiento) conforme especificación. & \cbox \\
		Registro de profundidades, diámetros, materiales y fotografías “as-built”. & \cbox \\
		Programa de \textbf{mantenimiento} (retiro de sedimentos en pretratamiento) definido. & \cbox \\
		\bottomrule
	\end{tabularx}
	
	
	
	%%%%%%%%%%%%%%%%%%%%%%%%%%%%%%%%%%%%%%%%%%%%%%%%%%
	
	
	
	
	\section{Construcción de plantas de tratamiento de aguas residuales (PTAR)}
	
	\subsection*{Descripción general}
	La \textbf{planta de tratamiento de aguas residuales (PTAR)} es una instalación hidráulica conformada por unidades de proceso destinadas a depurar el agua mediante operaciones físicas, químicas y biológicas.  
	Su correcta construcción garantiza la eficiencia del sistema, la durabilidad de las estructuras y la seguridad ambiental.  
	Una PTAR puede incluir: rejas, desarenadores, sedimentadores, tanques de aireación, filtros, digestores, clarificadores, pozos, tuberías, cámaras de distribución y obras de descarga.
	

	
	% =====================================================================
	\subsection{Checklist de supervisión}
	
	\subsubsection*{1. Replanteo, excavaciones y control del terreno}
	\begin{tabularx}{\linewidth}{>{\raggedright\arraybackslash}p{0.93\linewidth} >{\centering\arraybackslash}m{0.07\linewidth}}
		\toprule
		\textbf{Actividad a revisar / ordenar} & \textbf{Estado}\\
		\midrule
		Trazos y cotas de cada unidad conforme planos topográficos y estructurales. & \cbox \\
		Excavaciones con taludes estables o entibaciones adecuadas. & \cbox \\
		Control del nivel freático y drenaje temporal operativo. & \cbox \\
		Compactación del fondo y verificación de densidad del terreno base. & \cbox \\
		\bottomrule
	\end{tabularx}
	
	\subsubsection*{2. Cimentaciones y losas de fondo}
	\begin{tabularx}{\linewidth}{>{\raggedright\arraybackslash}p{0.93\linewidth} >{\centering\arraybackslash}m{0.07\linewidth}}
		\toprule
		\textbf{Actividad a revisar / ordenar} & \textbf{Estado}\\
		\midrule
		Colocación de plantilla o base de grava nivelada. & \cbox \\
		Refuerzo colocado según planos: diámetros, separaciones y recubrimientos. & \cbox \\
		Concreto con f’c y slump especificado; vibrado y curado correctos. & \cbox \\
		Juntas de construcción selladas y tratamiento impermeable aplicado. & \cbox \\
		\bottomrule
	\end{tabularx}
	
	\subsubsection*{3. Muros, tabiques y compartimentos}
	\begin{tabularx}{\linewidth}{>{\raggedright\arraybackslash}p{0.93\linewidth} >{\centering\arraybackslash}m{0.07\linewidth}}
		\toprule
		\textbf{Actividad a revisar / ordenar} & \textbf{Estado}\\
		\midrule
		Encofrado alineado y arriostrado; secciones conforme planos. & \cbox \\
		Acero de refuerzo vertical y horizontal colocado correctamente. & \cbox \\
		Tratamiento de juntas frías con aditivo o banda hidroexpansiva. & \cbox \\
		Curado del concreto y aplicación de sellador impermeabilizante. & \cbox \\
		\bottomrule
	\end{tabularx}
	
	\subsubsection*{4. Instalaciones hidráulicas y electromecánicas}
	\begin{tabularx}{\linewidth}{>{\raggedright\arraybackslash}p{0.93\linewidth} >{\centering\arraybackslash}m{0.07\linewidth}}
		\toprule
		\textbf{Actividad a revisar / ordenar} & \textbf{Estado}\\
		\midrule
		Colocación de tuberías, válvulas, accesorios y conexiones según planos. & \cbox \\
		Sellado estanco de pasos de tubería en muros o losas. & \cbox \\
		Montaje de bombas, sopladores, difusores y rejillas metálicas. & \cbox \\
		Cableado eléctrico y tableros protegidos contra humedad. & \cbox \\
		\bottomrule
	\end{tabularx}
	
	\subsubsection*{5. Impermeabilización y acabados}
	\begin{tabularx}{\linewidth}{>{\raggedright\arraybackslash}p{0.93\linewidth} >{\centering\arraybackslash}m{0.07\linewidth}}
		\toprule
		\textbf{Actividad a revisar / ordenar} & \textbf{Estado}\\
		\midrule
		Superficies pulidas, sin nidos ni grietas; sellado de juntas completo. & \cbox \\
		Aplicación de recubrimiento impermeable en tanques y muros en contacto con agua. & \cbox \\
		Pendientes internas verificadas hacia desagües o reboses. & \cbox \\
		Protección contra rayos UV en tuberías y equipos expuestos. & \cbox \\
		\bottomrule
	\end{tabularx}
	
	\subsubsection*{6. Pruebas hidráulicas y eléctricas}
	\begin{tabularx}{\linewidth}{>{\raggedright\arraybackslash}p{0.93\linewidth} >{\centering\arraybackslash}m{0.07\linewidth}}
		\toprule
		\textbf{Actividad a revisar / ordenar} & \textbf{Estado}\\
		\midrule
		\textbf{Prueba de llenado y estanqueidad} de tanques sin fugas ni pérdidas de nivel. & \cbox \\
		Prueba de presión en líneas de impulsión y válvulas operativas. & \cbox \\
		Verificación de funcionamiento de bombas, aireadores y tableros eléctricos. & \cbox \\
		Medición de niveles, caudales y tiempos de retención. & \cbox \\
		\bottomrule
	\end{tabularx}
	
	\subsubsection*{7. Puesta en marcha y control final}
	\begin{tabularx}{\linewidth}{>{\raggedright\arraybackslash}p{0.93\linewidth} >{\centering\arraybackslash}m{0.07\linewidth}}
		\toprule
		\textbf{Actividad a revisar / ordenar} & \textbf{Estado}\\
		\midrule
		Tanques y estructuras limpias, sin residuos ni obstrucciones. & \cbox \\
		Pruebas iniciales con agua limpia realizadas satisfactoriamente. & \cbox \\
		Entrega de planos \textbf{as-built}, manuales y protocolo de operación. & \cbox \\
		Personal capacitado para operación y mantenimiento de la PTAR. & \cbox \\
		\bottomrule
	\end{tabularx}
	
	
	
	
	
	%%%%%%%%%%%%%%%%%%%%%%%%%%%%%%%%%%%%%%%%%%%%%%%%%%%%%%%%%%%%%%%%
	
	
		  \newpage
	% Segunda introducción
	\newpage
	\stepcounter{intro}
	\section*{\centering \theintro\ Descripción y proceso constructivo de agua potable}
	\addcontentsline{toc}{section}{\theintro\ Descripción y proceso constructivo de elementos de concreto}
	\setcounter{section}{0}
	
	
	
	
	
	%%%%%%%%%%%%%%%%%%%%%%%%%%%%%%%%%%%%%%%%%%%%%%%%%%
			
	\section{Instalación de línea de conducción superficial y enterrada}
	
	\subsection*{Descripción general}
	La \textbf{línea de conducción} es el sistema de tuberías encargado de transportar agua desde una fuente (nacimiento, tanque o estación de bombeo) hasta una planta de tratamiento o red de distribución.  
	Puede ser \textbf{superficial} (sobre terreno o estructuras) o \textbf{enterrada} (bajo suelo o caminos). Su correcta instalación requiere control de pendientes, alineamientos, anclajes, juntas, válvulas y pruebas hidráulicas para asegurar la continuidad y estanqueidad del sistema.
	

	
	% =====================================================================
	\subsection{Checklist de supervisión}
	
	\subsubsection*{1. Replanteo y preparación del terreno}
	\begin{tabularx}{\linewidth}{>{\raggedright\arraybackslash}p{0.93\linewidth} >{\centering\arraybackslash}m{0.07\linewidth}}
		\toprule
		\textbf{Actividad a revisar / ordenar} & \textbf{Estado}\\
		\midrule
		Trazos y ejes de la línea definidos conforme planos topográficos y de diseño. & \cbox \\
		Pendientes verificadas con nivel o estación total. & \cbox \\
		Área de trabajo limpia, libre de obstáculos y con accesos seguros. & \cbox \\
		\bottomrule
	\end{tabularx}
	
	\subsubsection*{2. Excavación y cama de apoyo (líneas enterradas)}
	\begin{tabularx}{\linewidth}{>{\raggedright\arraybackslash}p{0.93\linewidth} >{\centering\arraybackslash}m{0.07\linewidth}}
		\toprule
		\textbf{Actividad a revisar / ordenar} & \textbf{Estado}\\
		\midrule
		Zanja excavada con ancho y profundidad según especificaciones. & \cbox \\
		Fondo nivelado, sin material suelto ni agua acumulada. & \cbox \\
		Colocación de \textbf{cama de arena o grava fina} compactada y nivelada. & \cbox \\
		Protección de las paredes de la zanja (taludes o entibación) según seguridad. & \cbox \\
		\bottomrule
	\end{tabularx}
	
	\subsubsection*{3. Montaje de tubería y uniones}
	\begin{tabularx}{\linewidth}{>{\raggedright\arraybackslash}p{0.93\linewidth} >{\centering\arraybackslash}m{0.07\linewidth}}
		\toprule
		\textbf{Actividad a revisar / ordenar} & \textbf{Estado}\\
		\midrule
		Tipo de material de tubería (PVC, PEAD, acero, hierro dúctil) conforme especificación. & \cbox \\
		Uniones limpias y lubricadas; \textbf{alineación} correcta de campanas y espigas. & \cbox \\
		Verificación de juntas (empaques de goma o soldadura) sin fugas. & \cbox \\
		Colocación de válvulas, ventosas y accesorios en puntos indicados. & \cbox \\
		Protección de tuberías metálicas con recubrimiento o cinta anticorrosiva. & \cbox \\
		\bottomrule
	\end{tabularx}
	
	\subsubsection*{4. Soportes y anclajes (líneas superficiales)}
	\begin{tabularx}{\linewidth}{>{\raggedright\arraybackslash}p{0.93\linewidth} >{\centering\arraybackslash}m{0.07\linewidth}}
		\toprule
		\textbf{Actividad a revisar / ordenar} & \textbf{Estado}\\
		\midrule
		Bancadas o apoyos construidos según planos y espaciados uniformemente. & \cbox \\
		Alineación y nivel de la tubería verificados antes del colado de anclajes. & \cbox \\
		Instalación de \textbf{anclajes de empuje} en codos, tees o reducciones. & \cbox \\
		Aplicación de pintura protectora o recubrimiento anticorrosivo. & \cbox \\
		\bottomrule
	\end{tabularx}
	
	\subsubsection*{5. Pruebas hidráulicas y control de calidad}
	\begin{tabularx}{\linewidth}{>{\raggedright\arraybackslash}p{0.93\linewidth} >{\centering\arraybackslash}m{0.07\linewidth}}
		\toprule
		\textbf{Actividad a revisar / ordenar} & \textbf{Estado}\\
		\midrule
		\textbf{Prueba hidrostática} realizada según especificación (presión, tiempo y pérdidas). & \cbox \\
		\textbf{Prueba de estanqueidad} en uniones y válvulas sin fugas visibles. & \cbox \\
		Registro de presión, tiempo y comportamiento de manómetro. & \cbox \\
		Corrección de fugas o reemplazo de accesorios defectuosos. & \cbox \\
		\bottomrule
	\end{tabularx}
	
	\subsubsection*{6. Relleno y compactación (líneas enterradas)}
	\begin{tabularx}{\linewidth}{>{\raggedright\arraybackslash}p{0.93\linewidth} >{\centering\arraybackslash}m{0.07\linewidth}}
		\toprule
		\textbf{Actividad a revisar / ordenar} & \textbf{Estado}\\
		\midrule
		Relleno con material fino hasta 30 cm sobre la generatriz superior de la tubería. & \cbox \\
		Compactación por capas delgadas hasta alcanzar densidad especificada. & \cbox \\
		Relleno final con material seleccionado o del sitio, libre de piedras grandes. & \cbox \\
		Restitución del terreno o pavimento según condiciones originales. & \cbox \\
		\bottomrule
	\end{tabularx}
	
	\subsubsection*{7. Acabados, señalización y control final}
	\begin{tabularx}{\linewidth}{>{\raggedright\arraybackslash}p{0.93\linewidth} >{\centering\arraybackslash}m{0.07\linewidth}}
		\toprule
		\textbf{Actividad a revisar / ordenar} & \textbf{Estado}\\
		\midrule
		Colocación de hitos, mojones o cinta de advertencia (líneas enterradas). & \cbox \\
		Protección física o barandales en líneas superficiales expuestas. & \cbox \\
		Entrega de planos “\textbf{as-built}” con coordenadas, cotas y materiales usados. & \cbox \\
		Área limpia, sin escombros, y libre de obstrucciones. & \cbox \\
		\bottomrule
	\end{tabularx}
	
	
	
	

	%%%%%%%%%%%%%%%%%%%%%%%%%%%%%%%%%%%%%%%%%%%%%%%%%%%%%%%%
	
	
	
	\section{Instalación de líneas de distribución}
	
	\subsection*{Descripción general}
	Las \textbf{líneas de distribución} conforman la red que entrega el agua desde tanques o estaciones de bombeo hacia los usuarios finales, integrando \textbf{tuberías principales y secundarias}, válvulas de seccionamiento, hidrantes, ventosas, descargas, y \textbf{acometidas domiciliarias con micromedición}.  
	Su correcta ejecución requiere una adecuada \textbf{sectorización}, control de pendientes y presiones, uniones estancas, pruebas hidráulicas, \textbf{desinfección} y un relleno/compactación que garantice la durabilidad del sistema.
	

	% =====================================================================
	\subsection{Checklist de supervisión}
	
	\subsubsection*{1. Replanteo, sectorización y seguridad vial}
	\begin{tabularx}{\linewidth}{>{\raggedright\arraybackslash}p{0.93\linewidth} >{\centering\arraybackslash}m{0.07\linewidth}}
		\toprule
		\textbf{Actividad a revisar / ordenar} & \textbf{Estado}\\
		\midrule
		Ejes, cotas y pendientes trazados conforme planos; cruces con servicios existentes localizados. & \cbox \\
		\textbf{Sectorización} definida: ubicaciones de válvulas de corte, hidrantes, ventosas y descargas. & \cbox \\
		Plan de manejo de tránsito, señalización y EPP implementados en el frente de trabajo. & \cbox \\
		\bottomrule
	\end{tabularx}
	
	\subsubsection*{2. Excavación y cama de apoyo}
	\begin{tabularx}{\linewidth}{>{\raggedright\arraybackslash}p{0.93\linewidth} >{\centering\arraybackslash}m{0.07\linewidth}}
		\toprule
		\textbf{Actividad a revisar / ordenar} & \textbf{Estado}\\
		\midrule
		Zanja con ancho y \textbf{profundidad} especificados; fondo nivelado y sin agua. & \cbox \\
		Colocación de \textbf{cama de arena/grava fina} (espesor y nivel) sin elementos punzantes. & \cbox \\
		Taludes estables o entibación instalada donde aplique; accesos seguros para operadores. & \cbox \\
		\bottomrule
	\end{tabularx}
	
	\subsubsection*{3. Montaje de tuberías y uniones}
	\begin{tabularx}{\linewidth}{>{\raggedright\arraybackslash}p{0.93\linewidth} >{\centering\arraybackslash}m{0.07\linewidth}}
		\toprule
		\textbf{Actividad a revisar / ordenar} & \textbf{Estado}\\
		\midrule
		Material y diámetro de tubería conforme especificación (PVC/PEAD/HD/Acero). & \cbox \\
		\textbf{Uniones} ejecutadas según sistema: cementado, junta elástica, termofusión, bridas o roscas. & \cbox \\
		Alineación y pendiente verificadas; protección anticorrosiva en partes metálicas. & \cbox \\
		\bottomrule
	\end{tabularx}
	
	\subsubsection*{4. Válvulas, hidrantes, ventosas y anclajes}
	\begin{tabularx}{\linewidth}{>{\raggedright\arraybackslash}p{0.93\linewidth} >{\centering\arraybackslash}m{0.07\linewidth}}
		\toprule
		\textbf{Actividad a revisar / ordenar} & \textbf{Estado}\\
		\midrule
		Instalación de \textbf{válvulas de seccionamiento} en cámaras con marco y tapa. & \cbox \\
		\textbf{Hidrantes} ubicados y ensamblados según plano (altura, válvulas y drenes). & \cbox \\
		\textbf{Ventosas} tramos altos; descargas en puntos bajos con protección de descarga. & \cbox \\
		\textbf{Anclajes de empuje} (codos, tees, reducciones) con concreto y pernos según diseño. & \cbox \\
		\bottomrule
	\end{tabularx}
	
	\subsubsection*{5. Acometidas domiciliarias y micromedición}
	\begin{tabularx}{\linewidth}{>{\raggedright\arraybackslash}p{0.93\linewidth} >{\centering\arraybackslash}m{0.07\linewidth}}
		\toprule
		\textbf{Actividad a revisar / ordenar} & \textbf{Estado}\\
		\midrule
		Derivaciones ejecutadas con \textbf{saddle}, tee o perforación controlada (según norma). & \cbox \\
		\textbf{Caja de medidor} instalada (cota, orientación y protección); válvula domiciliaria accesible. & \cbox \\
		Prueba de estanqueidad de acometida y registro en planos \textit{as-built}. & \cbox \\
		\bottomrule
	\end{tabularx}
	
	\subsubsection*{6. Pruebas hidráulicas, desinfección y muestreo}
	\begin{tabularx}{\linewidth}{>{\raggedright\arraybackslash}p{0.93\linewidth} >{\centering\arraybackslash}m{0.07\linewidth}}
		\toprule
		\textbf{Actividad a revisar / ordenar} & \textbf{Estado}\\
		\midrule
		\textbf{Prueba hidrostática}: presión/tiempo según especificación; sin caídas inadmisibles. & \cbox \\
		\textbf{Desinfección} (cloración): dosis, tiempo de contacto y \textbf{enjuague} controlados. & \cbox \\
		Muestreo de \textbf{calidad de agua} (cloro residual, turbidez, bacteriología) si corresponde. & \cbox \\
		\bottomrule
	\end{tabularx}
	
	\subsubsection*{7. Relleno, reposición y señalización}
	\begin{tabularx}{\linewidth}{>{\raggedright\arraybackslash}p{0.93\linewidth} >{\centering\arraybackslash}m{0.07\linewidth}}
		\toprule
		\textbf{Actividad a revisar / ordenar} & \textbf{Estado}\\
		\midrule
		Relleno inicial con material fino hasta 30 cm sobre la generatriz; \textbf{compactación} por capas. & \cbox \\
		Reposición de pavimentos, banquetas y áreas verdes según especificación municipal. & \cbox \\
		Señalización permanente (hitos, tapas, marcadores) y \textbf{planos as-built} actualizados. & \cbox \\
		\bottomrule
	\end{tabularx}
	

	
	
	
	

		



	%%%%%%%%%%%%%%%%%%%%%%%%%%%%%%%%%%%%%%%%%%%%%%%
	
	
	
	
	
	
	
	\section{Construcción de tanque de distribución}
	
	\subsection*{Descripción general}
	El \textbf{tanque de distribución} es una estructura hidráulica destinada a almacenar agua tratada y distribuirla hacia diferentes sectores de una red.  
	Generalmente se construye en concreto armado o mampostería reforzada, con cámara de válvulas, ventilación, rebose, drenaje y sistema de medición de nivel.  
	Su ejecución debe garantizar \textbf{hermeticidad, estabilidad estructural y correcta conexión} con las líneas de entrada y salida.
	
		
	% =====================================================================
	\subsection{Checklist de supervisión}
	
	\subsubsection*{1. Replanteo y excavación}
	\begin{tabularx}{\linewidth}{>{\raggedright\arraybackslash}p{0.93\linewidth} >{\centering\arraybackslash}m{0.07\linewidth}}
		\toprule
		\textbf{Actividad a revisar / ordenar} & \textbf{Estado}\\
		\midrule
		Ubicación y cotas del tanque conforme planos topográficos y estructurales. & \cbox \\
		Excavación con taludes o entibación según profundidad; fondo limpio y nivelado. & \cbox \\
		Verificación de capacidad portante del terreno y control de nivel freático. & \cbox \\
		\bottomrule
	\end{tabularx}
	
	\subsubsection*{2. Cimentación y losa de fondo}
	\begin{tabularx}{\linewidth}{>{\raggedright\arraybackslash}p{0.93\linewidth} >{\centering\arraybackslash}m{0.07\linewidth}}
		\toprule
		\textbf{Actividad a revisar / ordenar} & \textbf{Estado}\\
		\midrule
		Colocación de plantilla de concreto pobre o capa de nivelación. & \cbox \\
		Armado de refuerzo inferior según plano; verificación de recubrimiento. & \cbox \\
		Vaciado de concreto con f’c especificado; vibrado y curado adecuado. & \cbox \\
		Instalación de juntas de construcción con banda hidroexpansiva. & \cbox \\
		\bottomrule
	\end{tabularx}
	
	\subsubsection*{3. Muros del tanque}
	\begin{tabularx}{\linewidth}{>{\raggedright\arraybackslash}p{0.93\linewidth} >{\centering\arraybackslash}m{0.07\linewidth}}
		\toprule
		\textbf{Actividad a revisar / ordenar} & \textbf{Estado}\\
		\midrule
		Encofrado alineado y arriostrado; verificación de verticalidad y espesor. & \cbox \\
		Colocación de acero principal y estribos conforme diseño. & \cbox \\
		Tratamiento de juntas frías con sellador o aditivo impermeable. & \cbox \\
		Curado continuo del concreto por mínimo 7 días. & \cbox \\
		\bottomrule
	\end{tabularx}
	
	\subsubsection*{4. Losa de cubierta y acceso}
	\begin{tabularx}{\linewidth}{>{\raggedright\arraybackslash}p{0.93\linewidth} >{\centering\arraybackslash}m{0.07\linewidth}}
		\toprule
		\textbf{Actividad a revisar / ordenar} & \textbf{Estado}\\
		\midrule
		Armado y encofrado de la losa superior conforme planos estructurales. & \cbox \\
		Aperturas y \textbf{tapas de acceso} con marco metálico colocadas en posición. & \cbox \\
		\textbf{Ventilación y rebose} instalados antes del colado final. & \cbox \\
		Curado adecuado y acabado uniforme de la losa. & \cbox \\
		\bottomrule
	\end{tabularx}
	
	\subsubsection*{5. Instalaciones hidráulicas y válvulas}
	\begin{tabularx}{\linewidth}{>{\raggedright\arraybackslash}p{0.93\linewidth} >{\centering\arraybackslash}m{0.07\linewidth}}
		\toprule
		\textbf{Actividad a revisar / ordenar} & \textbf{Estado}\\
		\midrule
		Colocación de tuberías de entrada, salida, rebose y drenaje con sus válvulas respectivas. & \cbox \\
		Anclajes y soportes ejecutados según planos; sellos impermeables en pasos de muro. & \cbox \\
		Instalación de \textbf{cámara de válvulas} accesible y segura. & \cbox \\
		Verificación de operación manual y cierre hermético de válvulas. & \cbox \\
		\bottomrule
	\end{tabularx}
	
	\subsubsection*{6. Impermeabilización y acabados}
	\begin{tabularx}{\linewidth}{>{\raggedright\arraybackslash}p{0.93\linewidth} >{\centering\arraybackslash}m{0.07\linewidth}}
		\toprule
		\textbf{Actividad a revisar / ordenar} & \textbf{Estado}\\
		\midrule
		Aplicación de recubrimiento impermeable o aditivo hidrófugo en muros y losa. & \cbox \\
		Superficies sin grietas, nidos o desprendimientos; acabado liso. & \cbox \\
		Pendientes internas dirigidas hacia drenaje o válvula de limpieza. & \cbox \\
		Protección del concreto expuesto con pintura o sellador externo. & \cbox \\
		\bottomrule
	\end{tabularx}
	
	\subsubsection*{7. Pruebas de estanqueidad y control final}
	\begin{tabularx}{\linewidth}{>{\raggedright\arraybackslash}p{0.93\linewidth} >{\centering\arraybackslash}m{0.07\linewidth}}
		\toprule
		\textbf{Actividad a revisar / ordenar} & \textbf{Estado}\\
		\midrule
		\textbf{Prueba de llenado}: llenar el tanque y medir pérdidas de nivel en 24 y 48 horas. & \cbox \\
		Verificación visual de filtraciones o fugas en juntas y paredes. & \cbox \\
		Limpieza interior y cloración previa a la puesta en servicio. & \cbox \\
		Entrega de planos “\textbf{as-built}”, pruebas documentadas y registro fotográfico. & \cbox \\
		\bottomrule
	\end{tabularx}
	




%%%%%%%%%%%%%%%%%%%%%%%%%%%%%%%%%%%%%%%%%%%%%%%%%%%%%%%%%%








	\section{Construcción de tanques de captación de agua potable}
	
	\subsection*{Descripción general}
	Los \textbf{tanques de captación} son estructuras hidráulicas que permiten \textbf{almacenar, regular y distribuir} el agua proveniente de manantiales, nacimientos o sistemas de bombeo antes de su conducción hacia la planta de tratamiento o red de distribución.  
	Generalmente se construyen en concreto armado, mampostería reforzada o materiales prefabricados, con dispositivos de \textbf{entrada, rebose, limpieza, ventilación y medición de nivel}.  
	Su correcta ejecución es fundamental para asegurar la \textbf{hermeticidad}, durabilidad y calidad sanitaria del agua captada.
	

	
	% =====================================================================
	\subsection{Checklist de supervisión}
	
	\subsubsection*{1. Replanteo y excavación}
	\begin{tabularx}{\linewidth}{>{\raggedright\arraybackslash}p{0.93\linewidth} >{\centering\arraybackslash}m{0.07\linewidth}}
		\toprule
		\textbf{Actividad a revisar / ordenar} & \textbf{Estado}\\
		\midrule
		Ubicación conforme planos topográficos y estudio hidráulico. & \cbox \\
		Excavación con taludes o entibación según tipo de suelo y profundidad. & \cbox \\
		Control del nivel freático y drenaje temporal implementado. & \cbox \\
		Terreno base nivelado y con densidad de compactación especificada. & \cbox \\
		\bottomrule
	\end{tabularx}
	
	\subsubsection*{2. Cimentación y losa de fondo}
	\begin{tabularx}{\linewidth}{>{\raggedright\arraybackslash}p{0.93\linewidth} >{\centering\arraybackslash}m{0.07\linewidth}}
		\toprule
		\textbf{Actividad a revisar / ordenar} & \textbf{Estado}\\
		\midrule
		Colocación de plantilla de concreto pobre o capa de nivelación. & \cbox \\
		Armado de refuerzo principal, separadores y recubrimiento adecuados. & \cbox \\
		Colado del concreto f’c especificado con vibrado y curado continuo. & \cbox \\
		Instalación de juntas de construcción con sellador o banda hidroexpansiva. & \cbox \\
		\bottomrule
	\end{tabularx}
	
	\subsubsection*{3. Muros perimetrales y divisiones internas}
	\begin{tabularx}{\linewidth}{>{\raggedright\arraybackslash}p{0.93\linewidth} >{\centering\arraybackslash}m{0.07\linewidth}}
		\toprule
		\textbf{Actividad a revisar / ordenar} & \textbf{Estado}\\
		\midrule
		Encofrado firme, alineado y con los refuerzos verificados. & \cbox \\
		Concreto colado por capas delgadas; vibrado sin generar segregación. & \cbox \\
		Tratamiento de juntas frías y curado continuo durante 7 días mínimo. & \cbox \\
		Verificación de verticalidad, espesor y acabado liso de las paredes. & \cbox \\
		\bottomrule
	\end{tabularx}
	
	\subsubsection*{4. Losa de cubierta y acceso}
	\begin{tabularx}{\linewidth}{>{\raggedright\arraybackslash}p{0.93\linewidth} >{\centering\arraybackslash}m{0.07\linewidth}}
		\toprule
		\textbf{Actividad a revisar / ordenar} & \textbf{Estado}\\
		\midrule
		Encofrado, acero de refuerzo y pendientes hacia drenajes verificados. & \cbox \\
		Instalación de \textbf{tapas de acceso}, ventilaciones y tubos de rebose. & \cbox \\
		Colado del concreto con curado adecuado y acabado impermeable. & \cbox \\
		Revisión de espesores y juntas antes del desencofrado. & \cbox \\
		\bottomrule
	\end{tabularx}
	
	\subsubsection*{5. Instalación de tuberías y válvulas}
	\begin{tabularx}{\linewidth}{>{\raggedright\arraybackslash}p{0.93\linewidth} >{\centering\arraybackslash}m{0.07\linewidth}}
		\toprule
		\textbf{Actividad a revisar / ordenar} & \textbf{Estado}\\
		\midrule
		Colocación de tuberías de entrada, salida, rebose y drenaje conforme planos. & \cbox \\
		Sellado impermeable en pasos de muro o losa. & \cbox \\
		Instalación de válvulas de compuerta y purga con acceso para mantenimiento. & \cbox \\
		Revisión del anclaje de los tubos y protección contra corrosión. & \cbox \\
		\bottomrule
	\end{tabularx}
	
	\subsubsection*{6. Impermeabilización y acabados}
	\begin{tabularx}{\linewidth}{>{\raggedright\arraybackslash}p{0.93\linewidth} >{\centering\arraybackslash}m{0.07\linewidth}}
		\toprule
		\textbf{Actividad a revisar / ordenar} & \textbf{Estado}\\
		\midrule
		Aplicación de sellador o recubrimiento impermeabilizante certificado para agua potable. & \cbox \\
		Superficies internas pulidas, sin grietas ni poros abiertos. & \cbox \\
		Curado final con protección frente a la radiación solar. & \cbox \\
		\textbf{Protección exterior} con pintura o recubrimiento anticorrosivo. & \cbox \\
		\bottomrule
	\end{tabularx}
	
	\subsubsection*{7. Pruebas, limpieza y desinfección}
	\begin{tabularx}{\linewidth}{>{\raggedright\arraybackslash}p{0.93\linewidth} >{\centering\arraybackslash}m{0.07\linewidth}}
		\toprule
		\textbf{Actividad a revisar / ordenar} & \textbf{Estado}\\
		\midrule
		\textbf{Prueba de estanqueidad:} llenar y monitorear pérdidas de nivel durante 24–48 horas. & \cbox \\
		Limpieza interna con cepillado y enjuague con agua limpia. & \cbox \\
		\textbf{Desinfección} con solución de hipoclorito de sodio (cloro libre residual ≥ 0.5 mg/L). & \cbox \\
		Toma de muestra bacteriológica para validación sanitaria. & \cbox \\
		\bottomrule
	\end{tabularx}
	
	
	
	
	
	%%%%%%%%%%%%%%%%%%%%%%%%%%%%%%%%%%%%%%
	
	
	
	\section{Perforación de pozos de agua potable}
	
	\subsection*{Descripción general}
	La \textbf{perforación de pozos de agua potable} comprende las actividades para construir una captación subterránea segura: \emph{permisos y trazado}, \emph{perforación}, \emph{toma de muestras y registro litológico}, \emph{colocación de ademe y filtro}, \emph{empaque granular y sello sanitario}, \emph{desarrollo}, \emph{pruebas de bombeo} y \emph{desinfección}. La calidad del pozo depende del método de perforación, control de verticalidad, selección de filtros y del programa de desarrollo y pruebas.
	

	% =====================================================================
	\subsection{Checklist de supervisión}
	
	\subsubsection*{1. Replanteo, permisos y seguridad del sitio}
	\begin{tabularx}{\linewidth}{>{\raggedright\arraybackslash}p{0.93\linewidth} >{\centering\arraybackslash}m{0.07\linewidth}}
		\toprule
		\textbf{Actividad a revisar / ordenar} & \textbf{Estado}\\
		\midrule
		Permisos vigentes (ambientales, sanitarios, uso de agua) y servidumbres de acceso. & \cbox \\
		Ubicación georreferenciada (coordenadas) y cota del terreno confirmadas. & \cbox \\
		Área cercada y señalizada; manejo de residuos de lodo y recortes. & \cbox \\
		Plan de seguridad: EPP, bloqueo de equipos, control de derrames. & \cbox \\
		\bottomrule
	\end{tabularx}
	
	\subsubsection*{2. Perforación y control de calidad}
	\begin{tabularx}{\linewidth}{>{\raggedright\arraybackslash}p{0.93\linewidth} >{\centering\arraybackslash}m{0.07\linewidth}}
		\toprule
		\textbf{Actividad a revisar / ordenar} & \textbf{Estado}\\
		\midrule
		Método y sarta de perforación acordes a la geología (tricono/martillo DTH/barrena). & \cbox \\
		Propiedades del \textbf{lodo} (densidad, viscosidad) o uso de aire/espuma controlados. & \cbox \\
		Control de \textbf{verticalidad} y diámetro del hueco por intervalos. & \cbox \\
		Registro de nivel freático y caudales durante la perforación. & \cbox \\
		\bottomrule
	\end{tabularx}
	
	\subsubsection*{3. Muestras y registro litológico}
	\begin{tabularx}{\linewidth}{>{\raggedright\arraybackslash}p{0.93\linewidth} >{\centering\arraybackslash}m{0.07\linewidth}}
		\toprule
		\textbf{Actividad a revisar / ordenar} & \textbf{Estado}\\
		\midrule
		Toma de \textbf{cuttings/testigos} cada intervalo (p.ej., 1--3 m) correctamente rotulados. & \cbox \\
		Elaboración de \textbf{registro litológico} con profundidad, granulometría y acuíferos. & \cbox \\
		Identificación de zonas productivas para definir \textbf{longitud y ubicación} de filtros. & \cbox \\
		\bottomrule
	\end{tabularx}
	
	\subsubsection*{4. Revestimiento, filtros y empaque granular}
	\begin{tabularx}{\linewidth}{>{\raggedright\arraybackslash}p{0.93\linewidth} >{\centering\arraybackslash}m{0.07\linewidth}}
		\toprule
		\textbf{Actividad a revisar / ordenar} & \textbf{Estado}\\
		\midrule
		Tubos de \textbf{ademe} (material, diámetro, espesor) libres de deformaciones. & \cbox \\
		\textbf{Filtro ranurado} (ancho de ranura/slot) acorde a análisis granulométrico. & \cbox \\
		Colocación de \textbf{grava filtro} limpia y bien graduada en el espacio anular. & \cbox \\
		\bottomrule
	\end{tabularx}
	
	\subsubsection*{5. Sello sanitario y cabezal}
	\begin{tabularx}{\linewidth}{>{\raggedright\arraybackslash}p{0.93\linewidth} >{\centering\arraybackslash}m{0.07\linewidth}}
		\toprule
		\textbf{Actividad a revisar / ordenar} & \textbf{Estado}\\
		\midrule
		\textbf{Sello sanitario} (bentonita/cemento) sobre el empaque hasta superficie. & \cbox \\
		Cabezal sanitario con tapa hermética, \textbf{respiro} y puerto de muestreo. & \cbox \\
		Plataforma/brocal de concreto con pendiente y drenaje perimetral. & \cbox \\
		\bottomrule
	\end{tabularx}
	
	\subsubsection*{6. Desarrollo del pozo}
	\begin{tabularx}{\linewidth}{>{\raggedright\arraybackslash}p{0.93\linewidth} >{\centering\arraybackslash}m{0.07\linewidth}}
		\toprule
		\textbf{Actividad a revisar / ordenar} & \textbf{Estado}\\
		\midrule
		Método de desarrollo: \emph{air-lift}, pistoneo o bombeo progresivo. & \cbox \\
		Extracción de finos hasta alcanzar \textbf{claridad} (turbidez meta) estable. & \cbox \\
		Registro de niveles, caudal y arena producida durante el desarrollo. & \cbox \\
		\bottomrule
	\end{tabularx}
	
	\subsubsection*{7. Pruebas de bombeo y recuperación}
	\begin{tabularx}{\linewidth}{>{\raggedright\arraybackslash}p{0.93\linewidth} >{\centering\arraybackslash}m{0.07\linewidth}}
		\toprule
		\textbf{Actividad a revisar / ordenar} & \textbf{Estado}\\
		\midrule
		\emph{Step test}: escalones de caudal con lectura de abatimiento vs.\ tiempo. & \cbox \\
		\emph{Prueba a caudal constante}: duración (p.ej., 24 h) y monitoreo de niveles. & \cbox \\
		Curva de \textbf{recuperación} y estimación de parámetros hidráulicos. & \cbox \\
		\bottomrule
	\end{tabularx}
	
	\subsubsection*{8. Equipamiento, desinfección y entrega}
	\begin{tabularx}{\linewidth}{>{\raggedright\arraybackslash}p{0.93\linewidth} >{\centering\arraybackslash}m{0.07\linewidth}}
		\toprule
		\textbf{Actividad a revisar / ordenar} & \textbf{Estado}\\
		\midrule
		Selección e instalación de \textbf{bomba}, columna y válvula check según curva de operación. & \cbox \\
		\textbf{Cloración} del pozo (cálculo de dosis, tiempo de contacto y enjuague). & \cbox \\
		Ensayos de \textbf{calidad de agua} inicial (físico-químico y microbiológico). & \cbox \\
		Planos \textit{as-built}, ficha técnica y registro fotográfico. & \cbox \\
		\bottomrule
	\end{tabularx}
	
	
	
	
	%%%%%%%%%%%%%%%%%%%%%%%%%%%%%%%%%%%%%%%%%%%%%
	
	
	\section{Cisternas de agua (enterradas o semienterradas)}
	
	\subsection*{Descripción general}
	Las \textbf{cisternas} son depósitos (generalmente de concreto reforzado o prefabricados) destinados al \textbf{almacenamiento de agua} para consumo o servicios, ubicados \textbf{enterrados o semienterrados} para proteger el recurso y estabilizar la temperatura. Deben garantizar \textbf{hermeticidad}, acceso seguro para mantenimiento, ventilación, rebose y drenaje de limpieza, así como conexiones de entrada/salida con válvulas y medidores.
	
	
	
	% =====================================================================
	\subsection{Checklist de supervisión}
	
	% === FORMATO DE TABLAS UNIFICADO ===
	\renewcommand{\arraystretch}{1.2}
	\small
	
	\subsubsection*{1. Replanteo, excavación y seguridad}
	\begin{tabularx}{\linewidth}{
			>{\raggedright\arraybackslash}X
			>{\centering\arraybackslash}p{0.12\linewidth}
		}
		\toprule
		\textbf{Actividad a revisar / ordenar} & \textbf{Estado}\\
		\midrule
		Ubicación, ejes y cotas conforme planos topográficos/arquitectónicos. & \cbox \\
		Excavación con taludes estables o entibación; \textbf{bombeo} si hay agua. & \cbox \\
		Plataforma base nivelada y con densidad de compactación especificada. & \cbox \\
		Control de accesos, señalización y EPP (espacio confinado). & \cbox \\
		\bottomrule
	\end{tabularx}
	
	\subsubsection*{2. Cimentación y losa de fondo}
	\begin{tabularx}{\linewidth}{
			>{\raggedright\arraybackslash}X
			>{\centering\arraybackslash}p{0.12\linewidth}
		}
		\toprule
		\textbf{Actividad a revisar / ordenar} & \textbf{Estado}\\
		\midrule
		Plantilla de concreto pobre o capa de nivelación ejecutada. & \cbox \\
		Refuerzo inferior colocado (diámetros, separaciones y \textbf{recubrimientos}). & \cbox \\
		Concreto con f’c y slump especificado; vibrado y curado correctos. & \cbox \\
		Juntas de construcción con banda hidroexpansiva o sellador. & \cbox \\
		\bottomrule
	\end{tabularx}
	
	\subsubsection*{3. Muros perimetrales}
	\begin{tabularx}{\linewidth}{
			>{\raggedright\arraybackslash}X
			>{\centering\arraybackslash}p{0.12\linewidth}
		}
		\toprule
		\textbf{Actividad a revisar / ordenar} & \textbf{Estado}\\
		\midrule
		Encofrado alineado y arriostrado; espesores y plomos verificados. & \cbox \\
		Acero vertical y horizontal conforme planos; separadores suficientes. & \cbox \\
		Tratamiento de \textbf{juntas frías} y paso de tuberías con sellos estancos. & \cbox \\
		Curado continuo (≥ 7 días) y acabado liso sin nidos. & \cbox \\
		\bottomrule
	\end{tabularx}
	
	\subsubsection*{4. Losa/tapa y accesos}
	\begin{tabularx}{\linewidth}{
			>{\raggedright\arraybackslash}X
			>{\centering\arraybackslash}p{0.12\linewidth}
		}
		\toprule
		\textbf{Actividad a revisar / ordenar} & \textbf{Estado}\\
		\midrule
		Armado y encofrado de losa superior; pendientes hacia registros si aplica. & \cbox \\
		\textbf{Boca de acceso} con marco y tapa hermética; \textbf{ventilación} instalada. & \cbox \\
		Instalación de rebose y drenaje de limpieza con válvula accesible. & \cbox \\
		Escalera fija o grapas internas con pasamanos y protección. & \cbox \\
		\bottomrule
	\end{tabularx}
	
	\subsubsection*{5. Instalaciones hidráulicas y equipamiento}
	\begin{tabularx}{\linewidth}{
			>{\raggedright\arraybackslash}X
			>{\centering\arraybackslash}p{0.12\linewidth}
		}
		\toprule
		\textbf{Actividad a revisar / ordenar} & \textbf{Estado}\\
		\midrule
		Tuberías de \textbf{entrada/salida}, rebose y drenaje conforme planos (diámetros/material). & \cbox \\
		\textbf{Sellos} en pasos de muro; válvulas de compuerta y check accesibles. & \cbox \\
		Líneas de succión a bomba con canastilla/filtro; medición de nivel (boya, transmisor). & \cbox \\
		Protección anticorrosiva en elementos metálicos expuestos. & \cbox \\
		\bottomrule
	\end{tabularx}
	
	\subsubsection*{6. Impermeabilización interna y protección externa}
	\begin{tabularx}{\linewidth}{
			>{\raggedright\arraybackslash}X
			>{\centering\arraybackslash}p{0.12\linewidth}
		}
		\toprule
		\textbf{Actividad a revisar / ordenar} & \textbf{Estado}\\
		\midrule
		Aplicación de recubrimiento \textbf{apto para agua potable} (epóxico/cementoso). & \cbox \\
		Sellado de juntas y grietas; superficie sin poros abiertos. & \cbox \\
		Impermeabilización/perimetral externa si hay \textbf{freático} y control de \textbf{flotación}. & \cbox \\
		Relleno lateral por capas, compactado, sin dañar muros ni tuberías. & \cbox \\
		\bottomrule
	\end{tabularx}
	
	\subsubsection*{7. Pruebas, limpieza y puesta en servicio}
	\begin{tabularx}{\linewidth}{
			>{\raggedright\arraybackslash}X
			>{\centering\arraybackslash}p{0.12\linewidth}
		}
		\toprule
		\textbf{Actividad a revisar / ordenar} & \textbf{Estado}\\
		\midrule
		\textbf{Prueba de estanqueidad}: llenado y observación 24–48 h (sin pérdidas). & \cbox \\
		Limpieza interna y \textbf{desinfección} (cloración y enjuague). & \cbox \\
		Verificación de operación de válvulas, rebose y drenaje. & \cbox \\
		Entrega de \textbf{planos as-built} y registro fotográfico. & \cbox \\
		\bottomrule
	\end{tabularx}
	
	
	
	%%%%%%%%%%%%%%%%%%%%%%%%%%%%%%%%%%%%%%%%%%%%%%%%%%
	
	
	
	
	\section{Construcción de Pavimentos }
	
	\subsection*{a) Replanteo de curvas}
	El \textbf{replanteo de curvas} garantiza que el trazo geométrico de la carretera se ejecute conforme a planos y rasantes proyectadas.  
	Las actividades principales incluyen:
	\begin{itemize}[leftmargin=*, itemsep=2pt]
		\item Ubicación de puntos de control y referencias (PI, PC, PT, estaciones intermedias).
		\item Verificación de radios, longitudes de curva y tangentes con estación total o GPS.
		\item Marcado en campo mediante estacas o pintura sobre el terreno.
		\item Ajuste de alineación horizontal y pendiente transversal conforme al diseño.
	\end{itemize}
	\textit{Supervisión:} comprobar precisión del trazo, continuidad, transición suave y correcta orientación de las tangentes.
	
	% ----------------------------------------------------
	\subsection*{b) Estabilización de base, sub-base, sub-rasante y rasante}
	Este proceso busca otorgar \textbf{resistencia, durabilidad y soporte uniforme} a la estructura del pavimento.  
	Las etapas comprenden:
	\begin{itemize}[leftmargin=*, itemsep=2pt]
		\item \textbf{Subrasante:} nivelación, compactación al 95 \% Proctor Modificado, control de humedad y densidad.
		\item \textbf{Sub-base:} colocación de material granular clasificado, extendido por capas (≤ 20 cm) y compactación.
		\item \textbf{Base:} material pétreo o estabilizado con cemento o emulsión; control de espesor, gradación y densidad.
		\item \textbf{Rasante:} verificación de cotas y pendientes transversales según planos.
	\end{itemize}
	\textit{Supervisión:} registro de ensayos de densidad, granulometría y control topográfico de cada capa.
	
	% ----------------------------------------------------
	\subsection*{c) Proceso constructivo de pavimento flexible}
	El \textbf{pavimento flexible} está compuesto por capas bituminosas que distribuyen las cargas hacia las capas inferiores.  
	Etapas típicas:
	\begin{enumerate}[leftmargin=*, itemsep=2pt]
		\item Imprimación de la base con emulsión asfáltica.
		\item Colocación de la mezcla asfáltica en caliente (mezcla densa o SMA).
		\item Compactación con rodillos estáticos y vibratorios hasta alcanzar la densidad de diseño.
		\item Aplicación de sello o riego de niebla según especificación.
	\end{enumerate}
	\textit{Control:} temperatura de mezcla, contenido de ligante, espesor y densidad in situ.
	
	% ----------------------------------------------------
	\subsection*{d) Proceso constructivo de pavimento rígido}
	El \textbf{pavimento rígido} transmite las cargas a través de la losa de concreto hidráulico, minimizando deformaciones.  
	Pasos principales:
	\begin{enumerate}[leftmargin=*, itemsep=2pt]
		\item Preparación y nivelación de la base estabilizada.
		\item Colocación de moldes laterales, barras de transferencia y refuerzo (si aplica).
		\item Vaciado de concreto con f’c ≥ 280 kg/cm²; vibrado y acabado superficial.
		\item Corte de juntas de contracción, curado y sellado.
	\end{enumerate}
	\textit{Supervisión:} control de slump, dosificación, curado, alineación de juntas y acabado antideslizante.
	
	% ----------------------------------------------------
	\subsection*{e) Proceso constructivo de pavimento semi-rígido}
	El \textbf{pavimento semi-rígido} combina una base estabilizada con cemento y una capa asfáltica, logrando alta capacidad estructural y durabilidad.  
	Etapas:
	\begin{itemize}[leftmargin=*, itemsep=2pt]
		\item Mezclado y extendido de la base tratada con cemento (BTC).
		\item Curado de la BTC (mínimo 7 días) con control de fisuración.
		\item Aplicación de riego de liga y colocación de mezcla asfáltica de rodadura.
	\end{itemize}
	\textit{Control:} contenido de cemento, humedad óptima, curado y espesor final.
	
	% ----------------------------------------------------
	\subsection*{f) Proceso constructivo de afirmado mejorado}
	El \textbf{afirmado mejorado} se utiliza en caminos rurales y consiste en capas compactadas de material granular seleccionado con adición de estabilizantes.  
	Pasos generales:
	\begin{itemize}[leftmargin=*, itemsep=2pt]
		\item Preparación del terreno natural y escarificado superficial.
		\item Extendido del material granular (grava-arena) con control de espesor y humedad.
		\item Compactación en capas de 15–20 cm hasta alcanzar densidad especificada.
		\item Adición de estabilizantes (cal, cemento o emulsión) si lo requiere el diseño.
	\end{itemize}
	\textit{Supervisión:} verificar granulometría, contenido de finos, compactación y pendiente transversal para drenaje.
	
	
	
	
	%%%%%%%%%%%%%%%%%%%%%%%%%%%%%%%%%%%%%%%%
	
	
	\section{Instalación de artefactos, instalación de equipo eléctrico supervisión de acabados}
	
	% --- Formato uniforme de la tabla ---
	\small
	\renewcommand{\arraystretch}{1.2}
	
	\subsubsection*{Checklist de supervisión}
	\begin{tabularx}{\linewidth}{
			>{\raggedright\arraybackslash}X
			>{\centering\arraybackslash}p{0.12\linewidth}
		}
		\toprule
		\textbf{Actividad a revisar / ordenar} & \textbf{Estado} \\
		\midrule
		Instalación de piso y azulejo (trazo, nivelación, juntas, adhesivo y boquillas). & \cbox \\
		Repello, cernido y alisado de muros y cielos (espesores, plomo y acabado). & \cbox \\
		Instalación de bombas sumergibles y de succión (anclajes, válvulas, pruebas). & \cbox \\
		Instalación de artefactos de baño (sellos, conexiones AF/AC, pruebas de fugas). & \cbox \\
		Instalación de transformadores eléctricos (puesta a tierra, protecciones, pruebas). & \cbox \\
		Entubado eléctrico (conduit, cajas, fijaciones, radios de curvatura, continuidad). & \cbox \\
		Vibrado del concreto en elementos de acabado expuesto (compactación, sin nidos). & \cbox \\
		Instalación de techos metálicos (estructura secundaria, fijaciones y sellos). & \cbox \\
		\bottomrule
	\end{tabularx}
	
	
	
	
	%%%%%%%%%%%%%%%%%%%%%%%%%%%%%%%%%%%%%%%%%%%%%%%%%%%
	
	
	
	
	% ====== a) Instalación de piso y azulejo ======
	\subsubsection*{a) Instalación de piso y azulejo}
	\small\renewcommand{\arraystretch}{1.2}
	\begin{tabularx}{\linewidth}{
			>{\raggedright\arraybackslash}X >{\centering\arraybackslash}p{0.12\linewidth}}
		\toprule
		\textbf{Actividad a revisar / ordenar} & \textbf{Estado} \\
		\midrule
		Superficie limpia, nivelada y humedecida (si aplica). & \cbox \\
		Alineación y escuadra del trazo (ejes, maestras, cuerdas). & \cbox \\
		Adhesivo adecuado (tipo, tiempo abierto, espesor). & \cbox \\
		Colocación y nivelación de piezas (crucetas, mazo de goma). & \cbox \\
		Cortes en bordes y encuentros (zócalos, nichos, perímetros). & \cbox \\
		Pendientes en áreas húmedas y hacia desagües. & \cbox \\
		Boquilla/sellador aplicado y limpieza final. & \cbox \\
		\bottomrule
	\end{tabularx}
	
	% ====== b) Repello, cernido, alisado ======
	\subsubsection*{b) Repello, cernido y alisado}
	\small\renewcommand{\arraystretch}{1.2}
	\begin{tabularx}{\linewidth}{
			>{\raggedright\arraybackslash}X >{\centering\arraybackslash}p{0.12\linewidth}}
		\toprule
		\textbf{Actividad a revisar / ordenar} & \textbf{Estado} \\
		\midrule
		Preparación del sustrato (picado, limpieza, puente de adherencia). & \cbox \\
		Guías/maestras de nivel y plomo. & \cbox \\
		Dosificación del mortero y tamizado de arena. & \cbox \\
		Espesor por paños y tiempos entre capas. & \cbox \\
		Acabado con llana (cerrado/alisado) según especificación. & \cbox \\
		Juntas en encuentros y aristas vivas protegidas. & \cbox \\
		Curado y protección contra fisuras tempranas. & \cbox \\
		\bottomrule
	\end{tabularx}
	
	% ====== c) Bombas sumergibles y de succión ======
	\subsubsection*{c) Instalación de bombas sumergibles y de succión}
	\small\renewcommand{\arraystretch}{1.2}
	\begin{tabularx}{\linewidth}{
			>{\raggedright\arraybackslash}X >{\centering\arraybackslash}p{0.12\linewidth}}
		\toprule
		\textbf{Actividad a revisar / ordenar} & \textbf{Estado} \\
		\midrule
		Base/anclajes y accesibilidad para mantenimiento. & \cbox \\
		Tuberías de succión/impulsión y válvula check/retención. & \cbox \\
		Cebado (si aplica) y protección contra golpe de ariete. & \cbox \\
		Alimentación eléctrica, protecciones y tablero de control. & \cbox \\
		Flotadores/sensores de nivel; arranque/paro automático. & \cbox \\
		Puesta a tierra y compatibilidad de potencia/caudal. & \cbox \\
		Prueba de caudal, fugas y vibraciones. & \cbox \\
		\bottomrule
	\end{tabularx}
	
	% ====== d) Artefactos de baño ======
	\subsubsection*{d) Instalación de artefactos de baño}
	\small\renewcommand{\arraystretch}{1.2}
	\begin{tabularx}{\linewidth}{
			>{\raggedright\arraybackslash}X >{\centering\arraybackslash}p{0.12\linewidth}}
		\toprule
		\textbf{Actividad a revisar / ordenar} & \textbf{Estado} \\
		\midrule
		Ubicación y nivelación (plantillas, pernos, distancias sanitarias). & \cbox \\
		Conexiones AF/AC y flexibles certificados. & \cbox \\
		Sellos en pedestales/sanitarios (silicón, cera). & \cbox \\
		Trampas y ventilación sanitaria operativas. & \cbox \\
		Pruebas de fugas y descargas. & \cbox \\
		Accesorios (ducha, grifería, porta rollo, espejos) firmes y alineados. & \cbox \\
		\bottomrule
	\end{tabularx}
	
	% ====== e) Transformadores eléctricos ======
	\subsubsection*{e) Instalación de transformadores eléctricos}
	\small\renewcommand{\arraystretch}{1.2}
	\begin{tabularx}{\linewidth}{
			>{\raggedright\arraybackslash}X >{\centering\arraybackslash}p{0.12\linewidth}}
		\toprule
		\textbf{Actividad a revisar / ordenar} & \textbf{Estado} \\
		\midrule
		Plataforma/base y anclajes; ventilación/clearances. & \cbox \\
		Relación de transformación y tensiones nominales correctas. & \cbox \\
		Conexiones primario/secundario; apriete y polaridad. & \cbox \\
		Puesta a tierra (electrodo, conductor, unión). & \cbox \\
		Protecciones (fusibles, interruptor, relés). & \cbox \\
		Pruebas (megger, continuidad, relación de transformación). & \cbox \\
		Señalización y barreras de seguridad. & \cbox \\
		\bottomrule
	\end{tabularx}
	
	% ====== f) Entubado eléctrico ======
	\subsubsection*{f) Entubado eléctrico}
	\small\renewcommand{\arraystretch}{1.2}
	\begin{tabularx}{\linewidth}{
			>{\raggedright\arraybackslash}X >{\centering\arraybackslash}p{0.12\linewidth}}
		\toprule
		\textbf{Actividad a revisar / ordenar} & \textbf{Estado} \\
		\midrule
		Trazado de rutas y separación de servicios (fuerza/datos). & \cbox \\
		Tipo de conduit (PVC/EMT) y diámetro conforme conductores. & \cbox \\
		Curvaturas y radios mínimos; uso de cajas en cambios de dirección. & \cbox \\
		Fijaciones y sujeciones a distancias reglamentarias. & \cbox \\
		Sellos, reducción de ingreso de agua/polvo. & \cbox \\
		Guía pasacables y continuidad de ductos. & \cbox \\
		\bottomrule
	\end{tabularx}
	
	% ====== g) Vibrado del concreto ======
	\subsubsection*{g) Vibrado del concreto}
	\small\renewcommand{\arraystretch}{1.2}
	\begin{tabularx}{\linewidth}{
			>{\raggedright\arraybackslash}X >{\centering\arraybackslash}p{0.12\linewidth}}
		\toprule
		\textbf{Actividad a revisar / ordenar} & \textbf{Estado} \\
		\midrule
		Slump y ventana de colocación dentro de especificación. & \cbox \\
		Inserción del vibrador vertical, sin tocar armaduras ni moldes. & \cbox \\
		Tiempo/frecuencia adecuados; solape entre inserciones. & \cbox \\
		Espesores por capas y recorrido uniforme. & \cbox \\
		Eliminación de nidos y oquedades; acabado superficial. & \cbox \\
		Curado inmediato (membrana/agua) y protección. & \cbox \\
		\bottomrule
	\end{tabularx}
	
	% ====== h) Techos metálicos ======
	\subsubsection*{h) Instalación de techos metálicos}
	\small\renewcommand{\arraystretch}{1.2}
	\begin{tabularx}{\linewidth}{
			>{\raggedright\arraybackslash}X >{\centering\arraybackslash}p{0.12\linewidth}}
		\toprule
		\textbf{Actividad a revisar / ordenar} & \textbf{Estado} \\
		\midrule
		Estructura secundaria (correas) alineada y anclada. & \cbox \\
		Pendiente y dirección del traslape conforme ficha técnica. & \cbox \\
		Fijaciones (tornillos autotaladrantes con arandela) y separación. & \cbox \\
		Sellador en cumbreras, limahoyas y perímetros. & \cbox \\
		Aislamiento termoacústico (si aplica) y barrera de vapor. & \cbox \\
		Canaletas/bajantes instalados y probados. & \cbox \\
		Seguridad en altura (línea de vida, arnés). & \cbox \\
		\bottomrule
	\end{tabularx}
	
	
	
	

\end{document}
